\documentclass[11pt,letterpaper]{article}
\usepackage[T1]{fontenc}
\usepackage[utf8]{inputenc}
\usepackage{lmodern}
\usepackage[margin=1in,headheight=15pt]{geometry}
\usepackage{amsmath,amssymb,amsthm,mathtools,mathrsfs}
\usepackage{microtype,booktabs,longtable,tabularx,array,enumitem}
\usepackage[dvipsnames]{xcolor}
\usepackage{fancyhdr}
\usepackage[colorlinks=true,linkcolor=MidnightBlue,citecolor=MidnightBlue,urlcolor=MidnightBlue]{hyperref}
\usepackage[nameinlink,noabbrev]{cleveref}
\hypersetup{pdftitle={Surcomplex Polynomial Algebra: Order Geometry, Newton Profiles, and Scale-Resolved Stability},pdfauthor={Research exposition},pdfsubject={Surcomplex polynomials; arbitrary-rank Hahn fields; root geometry; perturbations; residue duality}}
\pagestyle{fancy}\fancyhf{}
\fancyhead[L]{\small Surcomplex polynomial algebra}
\fancyhead[R]{\small Order, scales, and multiplicity}
\fancyfoot[C]{\thepage}
\setlength{\parskip}{0.25em}
\setlength{\parindent}{1.25em}
\setlist{leftmargin=2em,itemsep=0.15em,topsep=0.3em}
\setcounter{tocdepth}{2}
\numberwithin{equation}{section}
\newtheorem{theorem}{Theorem}[section]
\newtheorem{proposition}[theorem]{Proposition}
\newtheorem{lemma}[theorem]{Lemma}
\newtheorem{corollary}[theorem]{Corollary}
\theoremstyle{definition}
\newtheorem{definition}[theorem]{Definition}
\newtheorem{example}[theorem]{Example}
\theoremstyle{remark}\newtheorem{remark}[theorem]{Remark}
\newcommand{\No}{\mathbf{No}}
\newcommand{\SC}{\mathbf{SC}}
\newcommand{\C}{\mathbb C}
\newcommand{\R}{\mathbb R}
\newcommand{\N}{\mathbb N}
\newcommand{\Z}{\mathbb Z}
\newcommand{\Q}{\mathbb Q}
\newcommand{\OO}{\mathcal O}
\newcommand{\mm}{\mathfrak m}
\newcommand{\HH}{\mathscr H_\Gamma}
\newcommand{\st}{\operatorname{st}}
\newcommand{\supp}{\operatorname{supp}}
\newcommand{\ord}{\operatorname{ord}}
\newcommand{\Res}{\operatorname{Res}}
\newcommand{\Disc}{\operatorname{Disc}}
\newcommand{\Tr}{\operatorname{Tr}}
\newcommand{\Nm}{\operatorname{Nm}}
\newcommand{\Hom}{\operatorname{Hom}}
\newcommand{\Spec}{\operatorname{Spec}}
\newcommand{\lc}{\operatorname{lc}}
\newcommand{\conv}{\operatorname{conv}}
\newcommand{\ini}{\operatorname{in}}
\newcommand{\id}{\operatorname{id}}
\newcommand{\Bge}[2]{B_{\ge}(#1,#2)}
\newcommand{\Bgt}[2]{B_{>}(#1,#2)}
\newcommand{\abs}[1]{\left|#1\right|}
\newcommand{\sh}{^{\sharp}}
\newcolumntype{L}{>{\raggedright\arraybackslash}X}
\newcolumntype{P}[1]{>{\raggedright\arraybackslash}p{#1}}
\begin{document}
\hypersetup{pageanchor=false}
\begin{titlepage}
\centering
\vspace*{0.7cm}
{\large\scshape A continuation of the supplied surcomplex manuscripts\par}
\vspace{1cm}
{\Huge\bfseries Surcomplex\par Polynomial Algebra\par}
\vspace{0.55cm}
{\LARGE Order Geometry, Newton Profiles,\par and Scale-Resolved Stability\par}
\vspace{0.85cm}
{\large September 21, 2026\par}
\vspace{0.8cm}
\begin{minipage}{0.94\textwidth}
\begin{abstract}
We develop finite-degree polynomial algebra over the surcomplex class
$\SC=\No[i]$, separating ordinary-looking order geometry from the
non-Archimedean geometry of its natural valuation. All proofs take place in
set-sized algebraically closed Hahn workspaces, and all degrees are ordinary
finite integers.

The order-geometric part proves surcomplex versions of root bounds,
Gauss--Lucas, Jensen's disk theorem, and the sharp Schoenberg second-moment
inequality, including its equality case. A precisely delimited real-closed-field
transfer gives a polynomial Rouch\'e theorem on circles with arbitrary surreal
centers and radii, without introducing fine-continuous contours.

The scale-sensitive part is organized by one initial-polynomial identity.
It determines root multiplicities in every valuation ball and residue
direction, proves Newton's polygon rule for arbitrary-rank ordered exponent
groups, and yields exact polynomial image balls. Differentiation of that
identity shows that a ball containing $k\ge1$ roots contains exactly $k-1$
critical points, counted with multiplicity. A finite root-cluster tree and an
explicit residual polynomial locate the intervening critical directions.

We prove support-controlled coprime factor lifting, an optimal
$1/n$ valuation-H\"older matching of perturbed root multisets, and
separation-sensitive simple-root stability with a second-order error estimate.
Finite quotient algebras, Hermite interpolation, resultants, and a
collision-stable residue pairing complete the algebraic picture. Examples
include nested critical scales, fractional and non-Archimedean exponents,
and a sharp collision threshold. Classical predecessors and dependencies on
the three supplied manuscripts are distinguished throughout.
\end{abstract}
\end{minipage}
\vfill
\begin{minipage}{0.94\textwidth}\small
\textbf{Status.} This article supplies proofs of the stated analogues and
support estimates. It does not claim that these are first occurrences in the
literature or solutions of a named open conjecture. Several central statements
are classical over real-closed or valued fields. The exposition and its exact
formulations have not been independently refereed or proof-assistant verified.
\end{minipage}
\end{titlepage}
\hypersetup{pageanchor=true}
\pagenumbering{roman}
{\small\tableofcontents}
\clearpage
\pagenumbering{arabic}

\section{The polynomial problem and its relation to the supplied work}\label{sec:intro}

\subsection{What changes when the coefficients are surcomplex?}
A surcomplex number is $x+iy$ with $x,y\in\No$ and $i^2=-1$.
The user's reference uses this terminology for the algebraically closed
class field obtained by complexifying the surreal numbers. The foundational
normal-form and real-closedness facts are part of established surreal theory
\cite{Conway,Gonshor}; the Hahn-field algebraic-closure theorem will provide our
set-sized version \cite{Poonen}.

For a polynomial of degree $n$, algebraic closedness already guarantees $n$
roots counted with multiplicity. That observation is important but does not
answer the geometric questions. How do roots lie relative to a disk whose
radius is infinite? How many roots occupy a given infinitesimal subcluster?
Which critical points sit between two such subclusters? How much coefficient
precision guarantees a one-to-one matching of roots, and when can a small
perturbation make distinct roots collide?

Two different structures answer these questions. The modulus
$|x+iy|=\sqrt{x^2+y^2}$ is surreal-valued and satisfies the ordinary triangle
inequality. The natural valuation records a leading Hahn exponent and
satisfies an ultrametric inequality. They must not be identified. In
particular, an order disk $|z-a|\le r$ is not the same object as the valuation
ball $v(z-a)\ge v(r)$.

The article treats the classical finite algebra as a foundation, then
establishes detailed results in both geometries. The main proofs are finite
algebraic arguments, except for the explicitly identified real-closed-field
transfer theorem and the well-founded support recursion for factor lifting.
Neither ordinary compactness of surreal circles nor convergence of Hahn
partial sums is used.

\subsection{The three supplied manuscripts}
The supplied \emph{Surcomplex Analysis} \cite{SourceAnalysis} distinguishes
fine-local analytic germs, Hahn-coherent functions on thickenings of ordinary
complex domains, and the much stronger condition of coherence at every
surreal scale. Its local preparation and moving-pole arguments already
connect analytic germs to finite cluster polynomials. Its all-scale rigidity
theorem singles out polynomials as the entire functions in that rigid category.
We use this as motivation, not as a prerequisite for the polynomial results.

The supplied \emph{Hartogs Coherence, Finite Maps, and Residue Duality}
\cite{SourceResearch} develops parameter-dependent division and finite
intersection algebras. In particular, its one-variable residue and trace
identities already include the collision-stable pairing that appears later
here. We give a purely polynomial proof of that identity and explicitly
credit the overlap.

The supplied \emph{Global Divisors and Cohomological Obstructions}
\cite{SourceGlobal} explains why an infinite family of individually legal
local clusters need not come from one coherent entire datum: the union of
its symmetric coefficient supports may fail to be well ordered. The finite
polynomial setting avoids that particular infinite gluing obstruction, but
retains the importance of symmetric rather than individually labeled root
data. The final comparison returns to this distinction.

The notation $t^\gamma=\omega^{-\gamma}$, set-sized Hahn workspaces, strong
summability, and ordinary standard parts follow these manuscripts. We do not
silently identify their coefficient sheaves with ordinary convergent
power-series rings, or assume their proposed novelty assessments as facts.

\subsection{A guide to the main conclusions}
The order-geometric results are in \cref{sec:geometry,sec:moment,sec:transfer}.
The central scale identity is \cref{thm:initialroots}; its consequences include
Newton root counts, exact image balls, and the critical-point conservation
law \cref{thm:criticalballs}. The root-matching and separated-root estimates
are \cref{thm:holder,thm:separated}. Coprime factor lifting is proved with
an explicit support bound in \cref{thm:hensel}. The quotient-algebra and
residue package culminates in \cref{thm:residuepairing,thm:traceresidue}.

A useful conceptual chain is
\[
\begin{gathered}
\text{finite factorization + coefficient reduction}\\
\Downarrow\\
\text{initial polynomial at a center and scale}\\
\Downarrow\\
\text{root counts, image balls, critical clusters, and perturbation stability}.
\end{gathered}
\]
The second-moment inequality belongs to a separate chain, through
orthogonal compression of a diagonal matrix. This separation keeps the
valuation from being used in a modulus argument or conversely.

\section{Fields, supports, and the two notions of size}\label{sec:foundations}

\subsection{Every finite problem has a set-sized workspace}
We use a class set theory admitting the usual treatment of $\No$, for example
NBG with suitable choice. Every polynomial, support, matrix, and recursion
in this article is set-sized. Polynomial degree always means an element of
$\N$, not an ordinal or an infinite surreal integer.

Let $\Gamma\subset\No$ be a set-sized divisible ordered additive subgroup.
Set
\[
 F_\Gamma=\R((t^\Gamma)),\qquad
 K_\Gamma=\C((t^\Gamma)),\qquad t^\gamma=\omega^{-\gamma}.
\]
Here a Hahn series is a formal expression
\[
 A=\sum_{\gamma\in S}A_\gamma t^\gamma
\]
with well-ordered nonzero support $S$. Conway normal forms embed these fields
in $\No$ and $\SC$, respectively. Real and imaginary parts show that
$K_\Gamma=F_\Gamma[i]$.

\begin{proposition}[Algebraically closed workspaces]\label{prop:workspace}
Every set of surcomplex constants is contained in some $K_\Gamma$ with
$\Gamma$ set-sized and divisible. Such $K_\Gamma$ is algebraically closed,
and its ordered real subfield $F_\Gamma$ is real closed. A polynomial already
split in $K_\Gamma$ has no additional roots after passage to a larger
surcomplex workspace.
\end{proposition}
\begin{proof}
Take the union of the supports of the given constants and then the divisible
hull of the additive subgroup generated by that union. These operations
produce sets. The Hahn-field algebraic-closure theorem applies because
$\Gamma$ is divisible and the residue field $\C$ is algebraically closed
\cite[Corollary 4]{Poonen}. An ordered field whose extension by $i$ is
algebraically closed is real closed, giving the second assertion.
Finally, if $P=a_n\prod_{j=1}^n(z-\alpha_j)$ already in $K_\Gamma[z]$, this
same identity holds in every extension field. Every root in an extension
must equal one of the $\alpha_j$.
\end{proof}
We usually write $F,K$ for one chosen pair. A theorem stated for all
surcomplex centers and scales means that each particular finite selection
is placed in such a workspace. It does not require one set-sized field to
contain every surreal number.

\subsection{Modulus and valuation}
For $z=x+iy\in K$ put
\[
 \bar z=x-iy,\qquad |z|=\sqrt{z\bar z}=\sqrt{x^2+y^2}\in F_{\ge0}.
\]
Positive square roots exist by real closedness. Algebraic
Cauchy--Schwarz gives
\[
 |zw|=|z||w|,\qquad |z+w|\le |z|+|w|.
\]
For completeness, the inequality
$(xu+yv)^2\le(x^2+y^2)(u^2+v^2)$ is the identity that the difference is
$(xv-yu)^2$. It yields the triangle inequality after squaring two
nonnegative sides. These arguments use no limits.

For $A\ne0$ define
\[
 v(A)=\min\supp(A),\qquad \lc(A)=A_{v(A)},\qquad v(0)=+\infty.
\]
Then
\begin{equation}\label{eq:valuation}
 v(AB)=v(A)+v(B),\qquad
 v(A+B)\ge\min\{v(A),v(B)\},\qquad v(|A|)=v(A).
\end{equation}
The last identity follows from the leading positive coefficient of
$A\bar A$. All nonzero complex constants have valuation zero, so in
particular $v(m)=0$ for every nonzero integer $m$. This residue-characteristic
zero fact is decisive for differentiation.

Write
\[
 \OO_K=\{A:v(A)\ge0\},\qquad \mm_K=\{A:v(A)>0\}.
\]
These are the finite elements and infinitesimals, respectively, relative
to the ordinary real constants. The coefficient at exponent zero defines
\[
 \st:\OO_K\longrightarrow\C,\qquad
 \OO_K/\mm_K\cong\C.
\]
Every finite $A$ is uniquely $c+\varepsilon$ with $c\in\C$ and
$\varepsilon\in\mm_K$. A nonzero element is infinitesimal exactly when
its modulus is less than $1/m$ for every positive ordinary integer $m$.

\begin{definition}[Order disks and valuation balls]\label{def:balls}
For $a\in K$, $r\in F_{>0}$, and $\rho\in\Gamma$, define
\[
 D(a,r)=\{z:|z-a|<r\},\qquad
 \overline D(a,r)=\{z:|z-a|\le r\},
\]
\[
 \Bge a\rho=\{z:v(z-a)\ge\rho\},\qquad
 \Bgt a\rho=\{z:v(z-a)>\rho\}.
\]
A residue-direction ball at this scale is
\[
 a+t^\rho(c+\mm_K),\qquad c\in\C.
\]
It is the open valuation ball $\Bgt{a+t^\rho c}{\rho}$.
\end{definition}
For example, $2t^\rho$ is in $\Bge0\rho$ but is not in
$\overline D(0,t^\rho)$. Likewise $|1+1|=2$, so the modulus itself is
not ultrametric. Terms such as ``open'' and ``closed'' attached to valuation
balls describe their defining inequalities; the results below do not use
fine-topological connectedness.

\subsection{The support lemma and its exact use}
A family of Hahn series is \emph{strongly summable} if its union of supports
is well ordered and each exponent occurs in only finitely many members.
For a well-ordered $E\subset\Gamma_{>0}$ put
\[
 E^*=\{e_1+\cdots+e_k:k\ge0,\ e_j\in E\},
\]
including the empty sum zero.

\begin{lemma}[Hahn--Neumann support calculus]\label{lem:support}
If $S,T$ are well-ordered subsets of an ordered abelian group, then $S+T$
is well ordered and each exponent has finitely many representations $s+t$.
If $E$ is well ordered and positive, then $E^*$ is well ordered and each
exponent has finitely many representations by nondecreasing finite words
in $E$. Consequently sums indexed by one element of $S$ and an arbitrary
finite word in $E$ are strongly summable when each word contributes only
finitely many terms.
\end{lemma}
These are the classical support facts \cite{Neumann}; they are also the
summability mechanism in the supplied manuscripts. To see the final
consequence, first separate an exponent into $s+m$, with $m\in E^*$.
There are finitely many such pairs, finitely many nondecreasing words for
each $m$, and finitely many orderings of each finite word. Thus both well
ordering and finite contribution multiplicity have been checked.

In particular, if $B$ has positive support, then
$(1+B)^{-1}=\sum_{k\ge0}(-B)^k$ is an exact Hahn identity. It need not be
a limit of partial sums in the full surreal fine topology. For example,
when the exponent group contains $\omega$, all remainders of
$\sum_{k\ge0}t^k$ after finitely many terms exceed the positive tolerance
$t^\omega$. Nor does $n\delta$ have to exceed an arbitrary positive group
element as $n$ grows. We never substitute either false cofinality claim
for \cref{lem:support}.

\section{Finite polynomial algebra and multiplicity}\label{sec:algebra}

\subsection{Factorization, division, and the squarefree part}
\begin{theorem}[The surcomplex fundamental theorem of algebra]\label{thm:fta}
Every nonconstant polynomial $P\in\SC[z]$ of degree $n$ admits a
factorization
\[
 P(z)=a_n\prod_{j=1}^n(z-\alpha_j),\qquad a_n\ne0.
\]
The multiset of $\alpha_j$ is unique. All roots lie in an algebraically
closed set-sized Hahn field containing the coefficients.
\end{theorem}
\begin{proof}
Use \cref{prop:workspace}, choose one root, and divide by its linear
factor. Induction on the finite degree gives the factorization. In a field,
if $P(\alpha)=0$, some factor in any displayed factorization vanishes at
$\alpha$. Canceling that factor and inducting proves multiset uniqueness.
\end{proof}
This is an application of established algebraic closedness, not a new
proof of the surreal normal-form theorem.

For nonzero $Q\in K[z]$, the leading-term division procedure gives unique
$S,R\in K[z]$ with $P=QS+R$ and $\deg R<\deg Q$. Each subtraction lowers
a nonnegative integer degree, so the procedure terminates. The Euclidean
algorithm therefore produces a monic greatest common divisor $G$ and
polynomials $A,B$ with $G=AP+BQ$. Every ideal of $K[z]$ is principal:
choose a nonzero member of least degree and divide every other member by
it. These are finite algebraic facts unaffected by the sizes of the
coefficients.

If the distinct roots of $P$ are $\xi$ with multiplicities $m_\xi$, then
\begin{equation}\label{eq:logderivative}
 \frac{P'(z)}{P(z)}=\sum_\xi\frac{m_\xi}{z-\xi},\qquad
 \gcd(P,P')=\prod_\xi(z-\xi)^{m_\xi-1}
\end{equation}
for monic $P$, with the gcd made monic. The product rule proves the first
identity. At $\xi$, writing $P=(z-\xi)^mU$ with $U(\xi)\ne0$ shows
$\ord_\xi P'=m-1$, because $m\ne0$ in characteristic zero. This proves
the second identity. Thus the squarefree part is $P/\gcd(P,P')$.

\subsection{Finite Hermite interpolation and local factors}
\begin{theorem}[Polynomial Chinese remainders and jets]\label{thm:crt}
Let $P=\prod_\xi(z-\xi)^{m_\xi}$ be monic of degree $n$. There is an
isomorphism
\begin{equation}\label{eq:crt}
 K[z]/(P)\ \cong\ \prod_\xi K[u]/(u^{m_\xi}),\qquad
 H\longmapsto
 \left(\sum_{j<m_\xi}\frac{H^{(j)}(\xi)}{j!}u^j\right)_\xi.
\end{equation}
In particular, arbitrary prescribed finite jets at these distinct points
are realized by a unique polynomial of degree less than $n$.
\end{theorem}
\begin{proof}
Finite Taylor expansion identifies the kernel of the map to the $\xi$
factor with $((z-\xi)^{m_\xi})$. Distinct factors are relatively prime.
For two coprime factors $A,B$, choose $U,V$ with $UA+VB=1$.
Then $VB$ is $1$ modulo $A$ and $0$ modulo $B$, and $UA$ has the
opposite residues. These two polynomials combine arbitrary residue
classes. Induction gives surjectivity for all the finitely many factors.
The kernel is their product ideal: divisibility by pairwise coprime
factors is equivalent to divisibility by their product. Unique division
by $P$ supplies the degree bound and uniqueness.
\end{proof}
This is a genuinely finite interpolation theorem. It has no uniform
support obstruction of the infinite type in \cite{SourceGlobal}; only
finitely many coefficient supports are being combined. Nevertheless,
its inverse can introduce inverse powers of small separations. Already
interpolation at $0,t^\rho$ uses the polynomial $z/t^\rho$ to give the
values $0,1$.

\subsection{Symmetric coordinates and Newton sums}
Write
\[
 P(z)=z^n+c_1z^{n-1}+\cdots+c_n,\qquad
 s_k=\sum_{j=1}^n\alpha_j^k.
\]
Expanding the finite product gives Vieta's formula
$c_k=(-1)^ke_k(\alpha_1,\ldots,\alpha_n)$. The Newton identities are
\begin{align}
 s_k+c_1s_{k-1}+\cdots+c_{k-1}s_1+kc_k&=0 &&(1\le k\le n),\
 s_k+c_1s_{k-1}+\cdots+c_ns_{k-n}&=0 &&(k>n).
\end{align}
One proof expands \eqref{eq:logderivative} at infinity as a formal Laurent
series in $z^{-1}$ and compares coefficients after multiplying by $P$.
Every coefficient comparison is finite. Alternatively the second identity
follows by substituting each root in $P$ and summing; the initial identities
follow successively from Vieta. No infinite analytic evaluation is involved.

For later use, a monic polynomial in $\OO_K[z]$ has all roots in
$\OO_K$. Indeed, if $v(\alpha)<0$, then in
$\alpha^n+\sum_{j<n}a_j\alpha^j$ the leading term has strictly smaller
valuation than every other term, so the sum cannot be zero. Conversely,
a monic polynomial with finite roots has finite coefficients by Vieta.
Reduction therefore gives the precise multiset identity
\begin{equation}\label{eq:reductionfactor}
 \st(P)(z)=\prod_{j=1}^n(z-\st(\alpha_j)).
\end{equation}
Multiple roots of the reduction can split at many different infinitesimal
scales. Reduction records their total multiplicity, not their separation.

\section{Order-valued root geometry}\label{sec:geometry}

\subsection{Root bounds at arbitrary surreal radii}
\begin{proposition}[Cauchy and a radial bound]\label{prop:rootbounds}
For $P(z)=a_nz^n+\cdots+a_0$, $a_n\ne0$, every root satisfies
\begin{equation}\label{eq:cauchybound}
 |\alpha|<1+\max_{j<n}|a_j/a_n|.
\end{equation}
If $a_0\ne0$, it also satisfies
\[
 |\alpha|>\frac{1}{1+\max_{1\le j\le n}|a_j/a_0|}.
\]
Writing $M=\max_{j<n}|a_j/a_n|^{1/(n-j)}$, one has $|\alpha|\le2M$;
when $M=0$ all roots are zero.
\end{proposition}
\begin{proof}
Divide by $a_n$ and let $A=\max_{j<n}|a_j/a_n|$. If $A=0$ the statement
is immediate. For $r=|\alpha|\ge1+A>1$, the equation $P(\alpha)=0$
and the finite geometric identity give
\[
 r^n\le A\sum_{j=0}^{n-1}r^j
   =A\frac{r^n-1}{r-1}\le r^n-1,
\]
a contradiction. Apply this bound to the reciprocal polynomial
$z^nP(1/z)/a_0$ to get the lower bound. Finally, if $r>2M>0$, then
\[
 \sum_{j<n}|a_j/a_n|r^{j-n}
 \le\sum_{j<n}(M/r)^{n-j}<\sum_{k=1}^n2^{-k}<1,
\]
so the leading term cannot cancel. The required positive roots defining
$M$ exist in the real-closed field.
\end{proof}
The bounds hold literally for infinite or infinitesimal coefficients.
No Archimedean estimate was used. Applying them to $P(a+rz)$ gives
corresponding bounds at any center and scale.

\subsection{Gauss--Lucas without a separation theorem}
For finitely many points $\alpha_j\in K$, their $F$-convex hull consists
of $\sum_j\theta_j\alpha_j$, where $\theta_j\in F_{\ge0}$ and
$\sum_j\theta_j=1$. Allowing surreal, rather than only ordinary real,
weights is essential.

\begin{theorem}[Surcomplex Gauss--Lucas]\label{thm:gausslucas}
Every root of $P'$ lies in the $F$-convex hull of the roots of $P$.
More explicitly, if $w$ is a critical point but not a root, then
\begin{equation}\label{eq:barycentric}
 w=\frac{\displaystyle\sum_\xi
      \frac{m_\xi\xi}{|w-\xi|^2}}
        {\displaystyle\sum_\xi\frac{m_\xi}{|w-\xi|^2}}.
\end{equation}
The same convex hull contains all zeros of every nonzero higher derivative.
\end{theorem}
\begin{proof}
At a nonroot critical point the logarithmic derivative vanishes. Since
$1/(w-\xi)=(\bar w-\bar\xi)/|w-\xi|^2$, conjugating that equality gives
\[
 \sum_\xi\frac{m_\xi(w-\xi)}{|w-\xi|^2}=0.
\]
Every weight is positive; their sum is nonzero and can be divided out,
yielding \eqref{eq:barycentric}. A critical point which is itself a root
is already in the hull. Repeating the argument for successive derivatives
and using convexity proves the final assertion.
\end{proof}
Thus any order-convex disk or half-plane containing the roots also
contains the critical points. The displayed proof is finite and algebraic;
it does not call on compactness of a convex hull or a topological
separation theorem. Its classical analogue is Gauss--Lucas; modern formal
proofs of the ordinary theorem and related root geometry are available
in \cite{Eberl}.

\subsection{Jensen disks for real-surreal coefficients}
\begin{theorem}[Surcomplex Jensen disk theorem]\label{thm:jensen}
Suppose $P\in F[z]$. Every nonreal critical point belongs to at least
one closed disk having a pair of conjugate nonreal roots as endpoints
of a diameter. Equivalently, for some root pair $a\pm ib$ with $b\ne0$,
that critical point $w$ satisfies $|w-a|\le |b|$.
\end{theorem}
\begin{proof}
Conjugation preserves the root multiset. A critical point that is a root
lies on the appropriate disk boundary, so assume $w=x+iy$ is not a root
and $y\ne0$. For a real root $r$,
\[
 \frac1y\operatorname{Im}\frac1{w-r}=-\frac1{|w-r|^2}<0.
\]
For a conjugate pair $a\pm ib$, a direct calculation gives
\begin{equation}\label{eq:jensencalc}
 \frac1y\operatorname{Im}\left(\frac1{w-a-ib}+\frac1{w-a+ib}\right)
 =\frac{-2\bigl((x-a)^2+y^2-b^2\bigr)}
 {|w-a-ib|^2|w-a+ib|^2}.
\end{equation}
If $w$ were outside every indicated closed disk, every contribution to
$\operatorname{Im}(P'/P)(w)/y$, with its positive multiplicity, would be
negative. Their finite sum could not vanish. This contradicts $P'(w)=0$.
\end{proof}
In particular, a real-rooted polynomial has only real critical points.
The assertion concerns $F$-valued disks; it is not a statement about
ordinary pictures obtained by discarding infinitesimals.

\section{A sharp second-moment inequality for critical points}\label{sec:moment}
The next result is the surcomplex version of the Schoenberg inequality.
The ordinary inequality is established mathematics, proved using
matrix differentiators and related spectral methods
\cite{Pereira,KushelTyaglov}. We include an algebraic proof over a general
real-closed field, including equality, so that no Hilbert-space completeness
assumption is hidden in the transfer.

\subsection{A matrix whose eigenvalues are the critical points}
Let $P$ be monic with roots $\alpha_1,\ldots,\alpha_n$, repeated according
to multiplicity, and $n\ge2$. Set
\[
 D=\operatorname{diag}(\alpha_1,\ldots,\alpha_n),\quad
 e=\frac1{\sqrt n}(1,\ldots,1)^T,\quad \Pi=I-ee^*.
\]
Here $*$ is conjugate transpose. Algebraic Gram--Schmidt over $F$ gives
a real $n\times(n-1)$ matrix $U$ with $U^*U=I$ and $UU^*=\Pi$.
Define $B=U^*DU$.

\begin{lemma}[Differentiating compression]\label{lem:compression}
The characteristic polynomial of $B$ is $P'(z)/n$.
\end{lemma}
\begin{proof}
The unitary matrix $(e\ U)$ changes $zI-D$ into a block matrix whose
lower-right block is $zI_{n-1}-B$. Its upper-left cofactor is consequently
$\det(zI_{n-1}-B)$. In the original coordinates this cofactor equals
$e^*\operatorname{adj}(zI-D)e$. Since $D$ is diagonal, it is
\[
 \frac1n\sum_{i=1}^n\prod_{j\ne i}(z-\alpha_j)=\frac1nP'(z).
\]
All expressions are polynomial identities, so no invertibility at a root
of $P$ is being assumed.
\end{proof}

\begin{lemma}[Finite-dimensional Schur inequality]\label{lem:schur}
If $A\in\operatorname{Mat}_m(K)$ has eigenvalues $\lambda_1,\ldots,\lambda_m$
with multiplicity, then
\[
 \sum_{j=1}^m|\lambda_j|^2\le\Tr(AA^*).
\]
Equality holds exactly when $A$ is normal, that is, $AA^*=A^*A$.
\end{lemma}
\begin{proof}
Algebraic closedness provides an eigenvector. Normalize it using the
positive square root of its squared norm, extend it to an orthonormal
basis by Gram--Schmidt, and apply induction to the remaining square
block. This gives a unitary upper triangularization. In that basis,
$\Tr(AA^*)$ is the sum of the squared moduli of all entries, while the
sum on the left is the sum for the diagonal entries. Their difference
is the sum of the squared moduli of the strict upper-triangular entries.
It vanishes exactly when those entries vanish.

A diagonal matrix is normal, and normality is invariant under unitary
conjugation. Conversely an upper-triangular normal matrix is diagonal:
equality of the $(1,1)$ entries of $TT^*$ and $T^*T$ forces every
strict upper entry in the first row to vanish, and induction handles
the remaining block. This proves the equality criterion as well.
\end{proof}

\subsection{The inequality and its equality case}
\begin{theorem}[Sharp surcomplex Schoenberg inequality]\label{thm:schoenberg}
Let $w_1,\ldots,w_{n-1}$ be the roots of $P'$, with multiplicity.
Then
\begin{equation}\label{eq:schoenberg}
 \sum_{j=1}^{n-1}|w_j|^2
 \le \frac{n-2}{n}\sum_{i=1}^n|\alpha_i|^2
    +\frac1{n^2}\left|\sum_{i=1}^n\alpha_i\right|^2.
\end{equation}
For the root centroid $c=n^{-1}\sum_i\alpha_i$, equivalently
\begin{equation}\label{eq:variance}
 \boxed{\displaystyle
 \sum_{j=1}^{n-1}|w_j-c|^2
 \le\frac{n-2}{n}\sum_{i=1}^n|\alpha_i-c|^2.}
\end{equation}
Equality holds if and only if all the roots lie on an affine $F$-line
in $K$. The one-point configuration is included as collinear.
\end{theorem}
\begin{proof}
Apply \cref{lem:schur} to the compression $B$ and use
\cref{lem:compression}. Cyclic invariance of the finite matrix trace gives
\begin{align*}
 \Tr(BB^*)&=\Tr(D\Pi D^*\Pi)\\
 &=\left(1-\frac2n\right)\sum_i|\alpha_i|^2
     +\frac1{n^2}\left|\sum_i\alpha_i\right|^2.
\end{align*}
This is \eqref{eq:schoenberg}. Vieta for $P'$ gives
$\sum_jw_j=(n-1)c$; replacing every root by its difference from $c$
gives \eqref{eq:variance}.

For equality, let $v=\Pi De$ and $u=\Pi D^*e=\bar v$; the last equality
uses that $\Pi$ is real. Regard the compression as an operator on
$e^\perp$. Expanding $\Pi=I-ee^*$ and using $DD^*=D^*D$ gives
\[
 BB^*-B^*B\quad\text{represented on }e^\perp\text{ by}\quad uu^*-vv^*.
\]
If $v=0$, all roots equal $c$. Otherwise equality of these two rank-one
matrices is equivalent to $u=\zeta v$ for some $|\zeta|=1$. To justify
this elementary fact, their common nonzero image is a one-dimensional
space, so $u=\zeta v$; substituting forces $\zeta\bar\zeta=1$.
The coordinates of $v$ are $(\alpha_i-c)/\sqrt n$. Therefore normality
is equivalent to
\[
 \overline{\alpha_i-c}=\zeta(\alpha_i-c)\qquad(1\le i\le n).
\]
Choose a nonzero $d=\alpha_j-c$. Dividing these equalities shows
$\overline{(\alpha_i-c)/d}=(\alpha_i-c)/d$, so every such ratio lies
in $F$. All roots lie on $c+dF$. Conversely this collinearity makes
$\bar v$ a unit-modulus multiple of $v$, hence makes $B$ normal.
\Cref{lem:schur} now gives exactly the asserted equality case.
\end{proof}
The coefficient $(n-2)/n$ cannot be improved, since nonconstant collinear
configurations attain equality. Both sides may be infinite, finite, or
infinitesimal surreal quantities; they are compared in the ordered field,
not after taking standard parts. As a quick strict example, the roots
$1,i,-1,-i$ of $z^4-1$ have critical points all zero, so the inequality
is strict. Multiplying every root by any nonzero surcomplex scale gives
an equally valid strict example at that scale.

\section{Polynomial Rouch\'e on order circles}\label{sec:transfer}

\subsection{Exactly what real-closed-field transfer permits}
Tarski's theorem says that first-order formulas in ordered-field
arithmetic admit quantifier elimination, and in particular that all
real-closed fields satisfy the same first-order sentences
\cite{TarskiSurvey}. For any \emph{fixed} degree bound, a polynomial with
complex coefficients can be represented by finitely many pairs of real
coordinates. Polynomial identities, circle equations, and squared-modulus
inequalities are then formulas in that language.

Root counts with multiplicity can also be encoded. Choose a finite list
of roots, impose the coefficient identities for a complete factorization,
and impose, by a finite disjunction, that exactly $k$ entries of the list
lie in the specified disk. This counts repeated entries rather than
only distinct points. Degrees below a chosen bound are treated by a
finite disjunction according to the last nonzero coefficient.

This is a theorem scheme, one sentence for each fixed finite degree
bound. It does not transfer unrestricted quantification over functions,
sequences, or integers as a distinguished subset. It also does not make
a proper class into a first-order set-sized model: transfer is applied
to $F_\Gamma$, after all parameters have been placed in that field.
Real-algebraic polynomial root counting by Sturm chains is developed
constructively in \cite{Eisermann}; the transfer below is compatible
with that algebraic viewpoint.

\begin{theorem}[Order-circle polynomial Rouch\'e]\label{thm:ordrouche}
Let $P,Q\in\SC[z]$ and let $a\in\SC$, $r\in\No_{>0}$. Suppose that
\begin{equation}\label{eq:ordrouche}
 |Q(z)-P(z)|<|P(z)|\qquad\text{for every }|z-a|=r.
\end{equation}
Then $P$ and $Q$ have the same number of roots in $D(a,r)$, counted
with multiplicity, and neither has a root on the circle.
\end{theorem}
\begin{proof}
Choose a real-closed workspace containing the coefficients, center,
and radius. Fix the actual degree bound. Over the ordinary complex
numbers, the polynomial case of Rouch\'e's theorem gives the asserted
implication. Its antecedent is the first-order formula
\[
 \forall x,y\;\bigl((x-a_R)^2+(y-a_I)^2=r^2
 \ \Longrightarrow\ |Q(x+iy)-P(x+iy)|^2<|P(x+iy)|^2\bigr).
\]
After expanding real and imaginary parts, every term is an ordered-field
polynomial. Encode the conclusion by complete factorization lists as
above. Real-closed-field transfer proves this sentence in the chosen
workspace. The full-surcomplex hypothesis implies its restriction to
that workspace, and \cref{prop:workspace} ensures its root lists already
contain every surcomplex root. Finally, a boundary root of $P$ contradicts
the strict inequality, and a boundary root of $Q$ would make its two
sides equal. Thus the boundary assertion also holds directly.
\end{proof}
This proof does not define a surreal contour integral. It is not the
same statement as the leading-coefficient Rouch\'e theorem on ordinary
thickened domains in \cite{SourceAnalysis}: its hypothesis concerns the
full order circle at the chosen surreal scale.

\begin{corollary}[Monomial dominance]\label{cor:pellet}
If $P(z)=\sum_{j=0}^na_jz^j$ and for some $k$ and $r>0$ one has
\[
 |a_k|r^k>\sum_{j\ne k}|a_j|r^j,
\]
then $P$ has exactly $k$ roots in $D(0,r)$ and none on its boundary.
\end{corollary}
\begin{proof}
On the circle the other monomials sum to a quantity of modulus strictly
less than $|a_kz^k|$. Apply \cref{thm:ordrouche} to $a_kz^k$ and $P$.
The monomial has exactly $k$ roots in the disk, including multiplicity
at zero; the constant case $k=0$ gives none.
\end{proof}
No strict inequality is required among the valuations of the terms.
The criterion can distinguish two terms at the same valuation by their
leading ordinary moduli. The valuation criterion below instead preserves
a complete initial polynomial, and answers a different set of questions.

\section{The initial polynomial at a center and a scale}\label{sec:initial}

\subsection{Weighted coefficient valuation}
Fix a nonzero $P\in K[z]$, a center $a\in K$, and a scale exponent
$\rho\in\Gamma$. Its finite Taylor expansion is
\[
 P(a+Y)=\sum_{j=0}^n b_jY^j,\qquad b_j=P^{(j)}(a)/j!.
\]
Define
\begin{equation}\label{eq:weighted}
 w_{a,\rho}(P)=\min_{b_j\ne0}\{v(b_j)+j\rho\}=:\lambda,
 \qquad
 I_{a,\rho}(P)(X)=\st\bigl(t^{-\lambda}P(a+t^\rho X)\bigr).
\end{equation}
Reduction in the second expression is coefficientwise. The normalized
polynomial has finite coefficients, at least one a unit, so its initial
polynomial $I_{a,\rho}(P)\in\C[X]$ is nonzero. Set $w_{a,\rho}(0)=+\infty$.
Let
\begin{equation}\label{eq:active}
 J_{a,\rho}(P)=\{j:v(b_j)+j\rho=\lambda\},\quad
 j_- =\min J_{a,\rho}(P),\quad j_+=\max J_{a,\rho}(P).
\end{equation}
Then $\ord_0 I_{a,\rho}(P)=j_-$ and $\deg I_{a,\rho}(P)=j_+$.

\begin{lemma}[Multiplicative weighted valuation]\label{lem:gaussvaluation}
For nonzero polynomials $P,Q$,
\[
 w_{a,\rho}(PQ)=w_{a,\rho}(P)+w_{a,\rho}(Q),\qquad
 I_{a,\rho}(PQ)=I_{a,\rho}(P)I_{a,\rho}(Q).
\]
Also $w_{a,\rho}(P+Q)\ge\min\{w_{a,\rho}(P),w_{a,\rho}(Q)\}$.
\end{lemma}
\begin{proof}
Normalize both polynomials as in \eqref{eq:weighted}. Their reductions
are nonzero and their product is nonzero in the domain $\C[X]$.
Consequently the normalized product has coefficient valuation exactly
zero, proving both multiplicative identities. The additive inequality
follows coefficient by coefficient before reduction.
\end{proof}
This weighted valuation is often called a Gauss valuation. In contrast
to the modulus on $K$, it is ultrametric. No supremum over a topological
circle is needed to define it.

\subsection{The exact root formula}
\begin{theorem}[Initial polynomial and all residue-direction root counts]
\label{thm:initialroots}
Write $P=a_n\prod_{i=1}^n(z-\alpha_i)$, with multiplicities. For every
$a,\rho$ there is a nonzero constant $C_{a,\rho}\in\C$ such that
\begin{equation}\label{eq:initialroots}
 \boxed{\displaystyle
 I_{a,\rho}(P)(X)=C_{a,\rho}
 \prod_{v(\alpha_i-a)\ge\rho}
 \left(X-\st\frac{\alpha_i-a}{t^\rho}\right).}
\end{equation}
Moreover
\begin{equation}\label{eq:profileproduct}
 w_{a,\rho}(P)=v(a_n)+\sum_{i=1}^n\min\{\rho,v(\alpha_i-a)\}.
\end{equation}
Thus, counting multiplicity,
\begin{align}
 \#\{\alpha_i\in\Bge a\rho\}&=j_+,\label{eq:closedcount}\\
 \#\{\alpha_i\in\Bgt a\rho\}&=j_-,\label{eq:opencount}\\
 \#\{i:v(\alpha_i-a)=\rho\}&=j_+-j_- .\label{eq:shellcount}
\end{align}
For each $c\in\C$, the number of roots in
$a+t^\rho(c+\mm_K)$ equals the multiplicity of $c$ as a root of
$I_{a,\rho}(P)$.
\end{theorem}
\begin{proof}
It suffices to understand a linear factor, since
\cref{lem:gaussvaluation} makes both the weighted valuation and the
initial polynomial multiplicative. If $v(\alpha-a)\ge\rho$, write
$\beta=(\alpha-a)/t^\rho\in\OO_K$. Then
\[
 a+t^\rho X-\alpha=t^\rho(X-\beta),
\]
whose weighted valuation is $\rho$ and whose normalized reduction is
$X-\st(\beta)$. If $v(\alpha-a)<\rho$, put
$\delta=v(\alpha-a)$ and $\ell=\lc(\alpha-a)$. Normalizing
$t^\rho X-(\alpha-a)$ by $t^{-\delta}$ gives a polynomial with
reduction the nonzero constant $-\ell$. Its weighted valuation is
$\delta$. The leading scalar $a_n$ contributes its leading complex
coefficient and its valuation. Multiplying these finitely many results
proves \eqref{eq:initialroots} and \eqref{eq:profileproduct}.

A finite normalized root reduces to zero exactly when its valuation
is positive, and reduces to $c$ exactly when it lies in $c+\mm_K$.
Degree, order at zero, and multiplicity at $c$ in
\eqref{eq:initialroots} therefore give all the asserted counts.
\end{proof}
The theorem holds for negative, positive, finite, and infinite surreal
exponents $\rho$. It packages all roots in the chosen ball by a single
ordinary complex polynomial. The roots outside that ball contribute
only a nonzero scalar to the reduction, but their valuations still
contribute to \eqref{eq:profileproduct}.

\subsection{A valuation Rouch\'e theorem}
\begin{corollary}[Preservation of the complete initial polynomial]
\label{cor:valrouche}
If $P\ne0$ and
\begin{equation}\label{eq:valrouche}
 w_{a,\rho}(Q-P)>w_{a,\rho}(P),
\end{equation}
then $w_{a,\rho}(Q)=w_{a,\rho}(P)$ and
$I_{a,\rho}(Q)=I_{a,\rho}(P)$. Consequently their root counts agree
in the closed ball, the open ball, and every residue-direction ball at
that center and scale.
\end{corollary}
\begin{proof}
Normalize by $t^{-w_{a,\rho}(P)}$ after replacing $z$ by $a+t^\rho X$.
The difference then has all coefficients of positive valuation, so
its reduction is zero and cannot cancel the nonzero reduction of $P$.
Apply \cref{thm:initialroots}.
\end{proof}
The degrees of $P$ and $Q$ need not agree. Additional roots of the
higher-degree polynomial may lie outside the selected ball. The strict
inequality is necessary for this general statement: $P=z$ and
$Q=z+1$ at scale zero have different residue directions even though the
error has the same, rather than a smaller, weighted size.

\section{Newton profiles over arbitrary ordered exponent groups}\label{sec:newton}

\subsection{The polygon without an Archimedean assumption}
At center zero put
\[
 \phi_P(\rho)=\min_j\{v(a_j)+j\rho\}.
\]
This is a minimum of finitely many affine expressions in the ordered
divisible group $\Gamma$. One can organize it by its finite breakpoints
without embedding $\Gamma$ in $\R$. Candidate breakpoints have the form
$(v(a_j)-v(a_k))/(k-j)$, which belongs to $\Gamma$ by divisibility.
Only those at which at least two indices attain the minimum matter.

\begin{theorem}[Newton's root-valuation rule]\label{thm:newton}
At a finite exponent $\rho$, the number of nonzero roots of valuation
$\rho$ is the difference between the largest and smallest indices
attaining $\phi_P(\rho)$. Every valuation of a nonzero root is such a
breakpoint. Equivalently, a lower Newton edge joining indices $j<k$
with slope $-\rho$ accounts for $k-j$ roots of valuation $\rho$.
Roots equal to zero contribute their exact multiplicity separately.
\end{theorem}
\begin{proof}
The count at a fixed scale is \eqref{eq:shellcount}. If a nonzero root
has valuation $\rho$, that count is positive, so at least two indices
attain the minimum. Conversely a positive index difference produces
exactly that many roots by the same formula. A lower edge is precisely
a set of coefficient points on the line $y+\rho x=\phi_P(\rho)$,
with all other coefficient points above it. Its horizontal span is
therefore the difference of the extreme active indices. This definition
uses order and rational scalar multiplication in $\Gamma$, not a real
Euclidean drawing.
\end{proof}
Equation \eqref{eq:profileproduct} gives an equivalent useful identity:
\[
 \phi_P(\rho)=v(a_n)+\sum_i\min\{\rho,v(\alpha_i)\}.
\]
Its slope, where one active index is fixed, counts roots above the
current scale. The familiar Newton polygon is therefore a finite
encoding of the root-valuation multiset. The standard rank-one analytic
version and valuation polygons are discussed in \cite[\S3]{Benedetto};
the proof above does not require a real-valued norm or completeness.

\subsection{A cubic with two distinct root scales}
\begin{example}\label{ex:newtoncubic}
Let $t=\omega^{-1}$ and
\[
 P(z)=z^3-t^2z-t^5.
\]
The nonzero coefficient points are $(0,5),(1,2),(3,0)$ and
\[
 \phi_P(\rho)=\min\{5,2+\rho,3\rho\}.
\]
At $\rho=1$ the active indices are $1,3$, so two roots have valuation
one. At $\rho=3$ the active indices are $0,1$, so one root has valuation
three. More precisely,
\[
 I_{0,1}(P)=X^3-X,\qquad I_{0,3}(P)=-X-1.
\]
Hence the two valuation-one roots have leading terms $t$ and $-t$,
whereas the remaining root has leading term $-t^3$.

The simple-root recursion at each reduced root gives, in ordinary
integer-power Hahn notation,
\begin{align*}
 \alpha_+&=t+\tfrac12t^3-\tfrac38t^5+O(t^7),\\
 \alpha_-&=-t+\tfrac12t^3+\tfrac38t^5+O(t^7),\\
 \alpha_0&=-t^3-t^7+O(t^{11}).
\end{align*}
For instance setting $z=tu$ gives $u^3-u-t^2=0$, whose reduced roots
$1,-1$ have derivative $2$. Setting $z=t^3u$ gives
$t^4u^3-u-1=0$, with reduced root $-1$ and derivative $-1$.
Here $O(t^k)$ means a Hahn power series with support in integers at
least $k$, not a numerical asymptotic assertion obtained by varying $t$.
\end{example}

\subsection{Infinitely ranked and nondiscrete inputs}
\begin{example}[A scale beyond all integer multiples]\label{ex:highrank}
In a divisible workspace containing $1$ and $\omega$, take
\[
 P(z)=(z-t)(z-t^\omega).
\]
Its two roots have valuations $1$ and $\omega$. The Newton profile has
both breakpoints and the preceding theorem applies unchanged. There is
no integer $N$ with $N>\omega$, and an iteration that only raises a
valuation through the sequence $1,2,3,\ldots$ cannot be advertised as
reaching the second scale. The finite initial-polynomial construction
has no such restriction.
\end{example}

\begin{example}[A coefficient not given by a discrete power series]
Let $\gamma_k=1-1/(k+2)$ and
$A=\sum_{k\ge0}t^{\gamma_k}$. Its support is well ordered of order type
$\omega$, even though all its exponents are below one. The polynomial
$z^2-A$ has two roots of valuation $1/4$, since $v(A)=1/2$.
Writing $A=t^{1/2}(1+B)$, the two roots are
\[
 \pm t^{1/4}\sum_{m\ge0}\binom{1/2}{m}B^m.
\]
The support of $B$ is well ordered and positive, so this binomial
identity is strongly summable by \cref{lem:support}. Substituting a fixed
ordinary real number $q\in(0,1)$ for $t$ in the original expression for
$A$ would not give a convergent series: $q^{\gamma_k}$ tends to $q$,
not zero. No such substitution is used in the Hahn construction.
\end{example}

\section{Polynomial image balls and finite mapping degrees}\label{sec:images}
The initial-polynomial identity also determines a polynomial map on
an entire valuation ball, not only its zero fiber.

\begin{theorem}[Exact image and fiber degree]\label{thm:imageballs}
Let $P$ be nonconstant and write
$P(a+Y)=b_0+\sum_{j=1}^nb_jY^j$. Put
\[
 \lambda=\min_{j\ge1}\{v(b_j)+j\rho\},\qquad
 d_-=\min\{j\ge1:v(b_j)+j\rho=\lambda\},\quad
 d_+=\max\{j\ge1:v(b_j)+j\rho=\lambda\}.
\]
Then
\begin{align}
 P(\Bge a\rho)&=\Bge{P(a)}{\lambda},\label{eq:imageclosed}\\
 P(\Bgt a\rho)&=\Bgt{P(a)}{\lambda}.\label{eq:imageopen}
\end{align}
Every target on the right of \eqref{eq:imageclosed} has exactly $d_+$
preimages in the indicated closed source ball, counted with polynomial
multiplicity. Every target on the right of \eqref{eq:imageopen} has
exactly $d_-$ preimages in the open source ball.
\end{theorem}
\begin{proof}
For $v(h)\ge\rho$, each term $b_jh^j$ has valuation at least $\lambda$;
for $v(h)>\rho$, each has valuation strictly greater than $\lambda$.
This proves both image inclusions.

Now let $v(y-P(a))\ge\lambda$. The polynomial
$t^{-\lambda}(P(a+t^\rho X)-y)$ has finite coefficients and its
reduction has degree $d_+$: the changed constant term cannot affect
the nonzero top coefficient. By \cref{thm:initialroots} it has exactly
$d_+$ roots in $\OO_K$. This proves surjectivity and the closed-fiber
count. If $v(y-P(a))>\lambda$, its reduced constant term is zero and
its order at zero is exactly $d_-$. The open-ball count from the same
theorem proves the remaining assertions.
\end{proof}
This is an arbitrary-rank polynomial counterpart of the usual
non-Archimedean disk-mapping theorem; compare
\cite[Corollary 3.11]{Benedetto}. The finite proof shows why neither a
complete real-valued norm nor a restricted-power-series convergence
hypothesis is needed here.

The two mapping degrees can differ. For $P(z)=z^2-z$ at $a=0,\rho=0$,
one has $\lambda=0$, $d_-=1$, $d_+=2$. Thus the finite elements map
onto the finite elements with degree two, while the infinitesimal
ball at zero maps bijectively to itself. The second preimage of an
infinitesimal target lies in the monad of one.

\begin{corollary}[Affine normalization of a finite polynomial problem]
\label{cor:normalization}
For finitely many polynomials of bounded finite degree, one can choose
one affine scale so that all their roots are finite in the rescaled
coordinate. Each nonzero polynomial can then be made monic with finite
coefficients by a scalar normalization.
\end{corollary}
\begin{proof}
Choose a workspace containing all their roots and a $\rho$ no larger
than the finitely many valuations of the nonzero root displacements
from the selected center. Then every $(\alpha-a)/t^\rho$ is finite.
Divide the rescaled polynomial by its leading coefficient and use
Vieta. The choice of a minimum among finitely many exponents is enough.
\end{proof}
This explains the naturalness of integral-coefficient hypotheses in the
perturbation theorems below. The normalization changes the numerical
valuation precision of the coefficients, so that precision must always
be measured after the stated change of coordinates.

\section{Critical points at every scale}\label{sec:critical}

\subsection{Differentiation commutes with nonconstant reduction}
\begin{lemma}[The derivative of an initial polynomial]\label{lem:initialderivative}
If $I=I_{a,\rho}(P)$ has degree $k\ge1$, then
\begin{equation}\label{eq:initialderivative}
 w_{a,\rho}(P')=w_{a,\rho}(P)-\rho,\qquad
 I_{a,\rho}(P')=I'.
\end{equation}
More generally, for $1\le r\le k$,
\[
 w_{a,\rho}(P^{(r)})=w_{a,\rho}(P)-r\rho,\qquad
 I_{a,\rho}(P^{(r)})=I^{(r)}.
\]
\end{lemma}
\begin{proof}
Differentiate the normalized polynomial
$t^{-\lambda}P(a+t^\rho X)$ with respect to $X$. It gives
$t^{\rho-\lambda}P'(a+t^\rho X)$. All coefficients remain finite, and
its reduction is $I'$. Because $\C$ has characteristic zero and $I$
is nonconstant, $I'$ is nonzero. Thus the displayed normalization has
weighted valuation exactly zero, proving \eqref{eq:initialderivative}.
For higher derivatives, $I^{(r)}\ne0$ whenever $r\le k$, and the same
argument applies directly.
\end{proof}
The hypothesis cannot simply be removed: a constant initial polynomial
has zero derivative and does not determine the next nonzero scale of
$P'$. Likewise this argument can fail in positive residue
characteristic, where a nonconstant reduction may have zero derivative.

\begin{theorem}[Exact critical-point conservation in occupied balls]
\label{thm:criticalballs}
Let $B$ be either $\Bge a\rho$ or $\Bgt a\rho$. If $B$ contains
exactly $k\ge1$ roots of $P$, counted with multiplicity, then it contains
exactly $k-1$ roots of $P'$, counted with multiplicity. More generally,
for $1\le r\le k$, it contains exactly $k-r$ roots of $P^{(r)}$.
\end{theorem}
\begin{proof}
For the closed ball, $\deg I_{a,\rho}(P)=k$. The previous lemma gives
an initial polynomial of degree $k-r$ for $P^{(r)}$, which counts its
closed-ball roots by \cref{thm:initialroots}.

For the open ball, $\ord_0 I_{a,\rho}(P)=k$. In particular its degree
is at least $k$, so the previous lemma still applies for $r\le k$.
In characteristic zero, $\ord_0 I^{(r)}=k-r$ in that range. This is
exactly the required open-ball root count.
\end{proof}
Every root is counted with its full multiplicity. For example, a root
of multiplicity $m$ contributes $m-1$ critical multiplicities at its
exact location before any intercluster critical point is considered.

The occupied-ball condition $k\ge1$ is essential. For $P=z^2-1$, the
infinitesimal ball at zero contains no root of $P$, but contains the
critical point zero. The theorem does not assert a negative count or
exclude such critical points.

\begin{corollary}[Ramification count on an arbitrary mapping ball]
\label{cor:ramificationball}
With the notation of \cref{thm:imageballs}, the closed source ball
contains exactly $d_+-1$ critical multiplicities of $P$ and the open
source ball contains exactly $d_--1$. In particular, the indicated
restriction of $P$ is bijective exactly when its mapping degree is one,
and this is equivalent to the absence of critical points in that ball.
\end{corollary}
\begin{proof}
Apply \cref{thm:criticalballs} to $P-P(a)$, whose closed and open
root counts are $d_+$ and $d_-$ by \cref{thm:imageballs}. Its derivative
is $P'$. A mapping degree one supplies one simple preimage for every
target. Conversely, if the degree is greater than one, choose a target
which is not one of the finitely many critical values. Such a target
exists in the infinite image ball; its degree-many preimages are
simple and distinct, so the map is not injective.
\end{proof}
The last qualification prevents confusing multiplicity one-to-one
counting with set-theoretic injectivity at a single exceptional fiber.

\subsection{The residual polynomial between child clusters}
\begin{theorem}[Critical residue directions inside a cluster]
\label{thm:directions}
Suppose the closed ball $\Bge a\rho$ contains $k$ roots and its
nonempty root residue directions are distinct $c_1,\ldots,c_s\in\C$
with multiplicities $m_1,\ldots,m_s$, where $\sum_jm_j=k$.
Each occupied child ball $a+t^\rho(c_j+\mm_K)$ contains $m_j-1$
critical multiplicities. The other critical points in the parent ball
have total multiplicity $s-1$, with residue directions given by the
roots, including multiplicity, of
\begin{equation}\label{eq:residualcritical}
 R(X)=\sum_{j=1}^s m_j\prod_{\ell\ne j}(X-c_\ell).
\end{equation}
No root of $R$ equals any $c_j$.
\end{theorem}
\begin{proof}
By \cref{thm:initialroots},
$I_{a,\rho}(P)=C\prod_j(X-c_j)^{m_j}$. Its derivative factors as
\[
 I' = C\prod_j(X-c_j)^{m_j-1}R(X).
\]
The leading coefficient of $R$ is $\sum_jm_j=k\ne0$, so its degree
is $s-1$. Also
$R(c_j)=m_j\prod_{\ell\ne j}(c_j-c_\ell)\ne0$.
Now apply \cref{lem:initialderivative,thm:initialroots} to read every
critical direction and its multiplicity from $I'$.
\end{proof}
For two child directions, the intervening critical direction is
\[
 \frac{m_1c_2+m_2c_1}{m_1+m_2}.
\]
The weights in this formula are the opposite cluster multiplicities;
it is generally not the centroid of the roots. For $X^m(X-1)$ the
nonzero critical direction is $m/(m+1)$, whereas the root centroid is
$1/(m+1)$. More generally the directions in \eqref{eq:residualcritical}
lie in the ordinary complex convex hull of the child directions by
Gauss--Lucas applied to $I$.

\subsection{A finite root-cluster tree}
For a finite set of distinct roots with their positive multiplicities,
construct a tree as follows. If there is more than one distinct root,
choose one root $a$ and let $\rho$ be the minimum valuation of a pairwise
difference in the set. The closed ball $\Bge a\rho$ contains the whole
set. Partition it into the equivalence classes
$v(\alpha-\beta)>\rho$. These are its occupied residue-direction
children. Recurse on each child with more than one distinct root.
A leaf is a single root location, retaining its multiplicity.

At each nonleaf node there are at least two children, because the
minimum separation valuation is attained. Recursion strictly reduces
the number of distinct roots in a child, so this is a finite tree even
when the valuations have arbitrarily high ordered-group rank. Each
internal node has an actual scale and a residual polynomial as in
\eqref{eq:residualcritical}.

\begin{corollary}[Critical allocation by the root tree]\label{cor:tree}
For a degree-$n$ polynomial with $d$ distinct roots, an internal node
with $s$ occupied children accounts for exactly $s-1$ critical
multiplicities in the part of its ball outside those children. A leaf
of root multiplicity $m$ accounts for $m-1$ at that root. These
allocations account for all $n-1$ critical multiplicities, since
\[
 \sum_{\text{internal nodes}}(s-1)=d-1,\qquad
 \sum_{\text{leaves}}(m-1)=n-d.
\]
\end{corollary}
\begin{proof}
At an internal node, \cref{thm:directions} gives $s-1$ outside its
occupied children. A child with $m$ total root multiplicities contains
$m-1$ critical multiplicities; its smallest closed root ball contains
that same number by \cref{thm:criticalballs}, so none are lost in the
intermediate part of the child ball. At a singleton leaf, the exact
root multiplicity gives $m-1$ at the root. The elementary tree identity
follows by counting edges: the sum of all child counts equals the
number of nodes minus one. If $d=1$, the direct multiplicity calculation
already gives the assertion. The top root ball contains all $n-1$
critical multiplicities, so the allocation is exhaustive.
\end{proof}
This is an allocation theorem, not a claim that pairwise root valuations
alone determine every critical point. The residual directions $c_j$
carry additional data, and further subclustering among roots of $R$
can require additional calculation.

\subsection{A three-scale quartic calculation}\label{subsec:nested}
\begin{example}[Critical points between nested roots]\label{ex:nested}
Consider
\[
 P(z)=z(z-t^3)(z-t)(z-1).
\]
Its roots form the nested clusters $\{0,t^3\}\subset\{0,t^3,t\}
\subset\{0,t^3,t,1\}$. The initial polynomials and critical directions
are as follows:
\begin{center}
\begin{tabular}{@{}ccll@{}}
\toprule
$\rho$ & $w_{0,\rho}(P)$ & $I_{0,\rho}(P)$ & $I_{0,\rho}(P')$\\
\midrule
$0$ & $0$ & $X^3(X-1)$ & $X^2(4X-3)$\\
$1$ & $3$ & $-X^2(X-1)$ & $-X(3X-2)$\\
$3$ & $7$ & $X(X-1)$ & $2X-1$\\
\bottomrule
\end{tabular}
\end{center}
Consequently the three critical points are simple and have leading terms
\begin{equation}\label{eq:nestedcritical}
 w_{\mathrm{outer}}=\tfrac34+\text{infinitesimal},\quad
 w_{\mathrm{middle}}=t\bigl(\tfrac23+\text{infinitesimal}\bigr),\quad
 w_{\mathrm{inner}}=t^3\bigl(\tfrac12+\text{infinitesimal}\bigr).
\end{equation}
At scale zero, two critical multiplicities remain in the zero monad
and one has direction $3/4$. At scale one, one remains in the smaller
zero child and one has direction $2/3$. At scale three, the final
critical point has direction $1/2$. The finite tree has three binary
splits, each supplying one intervening critical point.

The derivative can also be expanded directly:
\[
 P'(z)=4z^3-3(1+t+t^3)z^2+2(t+t^3+t^4)z-t^4.
\]
Its coefficient profile verifies the critical valuations $0,1,3$.
The residual calculation additionally supplies their leading complex
coefficients and shows why all three are simple.
\end{example}

\section{Coprime factor lifting with controlled Hahn support}\label{sec:hensel}

\subsection{A finite linear splitting}
Let $p_1,\ldots,p_s\in\C[z]$ be pairwise coprime monic polynomials of
degrees $d_i\ge1$, and let $n=\sum_id_i$, $P_0=\prod_ip_i$.
The map
\begin{equation}\label{eq:factorlinear}
 L:\bigoplus_{i=1}^s\C[z]_{<d_i}\longrightarrow\C[z]_{<n},\qquad
 (h_i)\longmapsto\sum_i h_i\prod_{j\ne i}p_j
\end{equation}
is an isomorphism. Indeed, reducing a zero image modulo $p_i$ forces
$h_i=0$, because the other factors are invertible modulo $p_i$; the
source and target have the same finite dimension. This is the
linearization of multiplication of monic factors, and the only linear
inverse used in the next theorem.

\begin{theorem}[Support-controlled coprime Hensel factorization]
\label{thm:hensel}
Let $P=P_0+E\in\OO_K[z]$ be monic of degree $n$, with $\deg E<n$
and all coefficients of $E$ infinitesimal. Let $S\subset\Gamma_{>0}$
be the union of their nonzero supports. There are unique monic
polynomials
\[
 P_i=p_i+A_i,\qquad \deg A_i<d_i,
\]
whose corrections have positive support and satisfy
\begin{equation}\label{eq:henselfactor}
 P=\prod_{i=1}^sP_i.
\end{equation}
All coefficients of all $A_i$ have support in $S^*\setminus\{0\}$.
The $P_i$ remain pairwise coprime over $K$.
\end{theorem}
\begin{proof}
Expanding the product, the equation is
\[
 E=L(A_1,\ldots,A_s)+N(A_1,\ldots,A_s),
\]
where every monomial in $N$ contains at least two correction factors.
Proceed by well-founded recursion on $S^*\setminus\{0\}$. At exponent
$\nu$, the nonlinear coefficient $N_\nu$ depends only on correction
coefficients at exponents strictly smaller than $\nu$. There are
finitely many contributions by \cref{lem:support}. Define
\begin{equation}\label{eq:henselrecursion}
 (A_{1,\nu},\ldots,A_{s,\nu})=L^{-1}(E_\nu-N_\nu).
\end{equation}
The degree bounds persist because both the linear splitting and every
product difference in the monic product have degree less than $n$.
The resulting Hahn coefficients have the asserted support and satisfy
the factorization coefficient by coefficient.

For uniqueness, take the least exponent at which two normalized factor
lists differ; the finite union of their supports is well ordered.
At that exponent the nonlinear difference is zero, since it contains
a smaller-exponent difference multiplied by a positive correction.
Injectivity of $L$ forces the alleged first difference to be zero,
a contradiction.

Finally, the coefficient matrix for the finite B\'ezout map of a pair
$P_i,P_j$ reduces to the corresponding invertible matrix for $p_i,p_j$.
Its determinant therefore has nonzero ordinary reduction and is a unit
of $\OO_K$. The adjugate gives an inverse and hence a B\'ezout identity
for $P_i,P_j$. They are coprime.
\end{proof}
No iteration through only the natural numbers is assumed to be cofinal
in the output support. Formula \eqref{eq:henselrecursion} is genuinely
a recursion on a well-ordered support monoid. It generalizes the
positive-support preparation mechanism in the supplied manuscripts to
an arbitrary finite coprime polynomial splitting.

\begin{corollary}[Canonical standard-part cluster factors]\label{cor:clusterfactor}
For monic $P\in\OO_K[z]$, write
$\st(P)=\prod_{c\in A}(z-c)^{m_c}$ with $A\subset\C$ finite.
Then $P$ has unique monic factors $P_c$ reducing to $(z-c)^{m_c}$.
The roots of $P_c$ are exactly the roots of $P$ in $c+\mm_K$.
\end{corollary}
\begin{proof}
Apply \cref{thm:hensel} to the pairwise coprime ordinary factors.
Every root of $P_c$ is finite and its reduction is a root of
$\st(P_c)$ by \eqref{eq:reductionfactor}, so all of them lie over $c$.
The product identity and degree counts prove that none are omitted.
\end{proof}
Individual roots of a multiple cluster need not have supports in the
original exponent group. For example $z^m-t$ lies in $K_\Z[z]$ but
splits as $\prod_{\zeta^m=1}(z-\zeta t^{1/m})$ after passing to a
divisible workspace. The theorem controls the \emph{coprime cluster
factors}, not a nonexistent simple-root splitting of their common
multiple ordinary reduction.

\subsection{Holomorphic parameters without shrinking the domain}
Let $U$ be an ordinary complex domain, and let
$\HH(U)=\OO(U)((t^\Gamma))$, with the convention and strong evaluation
of the supplied manuscripts.

\begin{theorem}[Parameter-dependent coprime polynomial factors]
\label{thm:parameterhensel}
Suppose $P_0(x,z)=\prod_ip_i(x,z)$, where the $p_i$ are ordinary
holomorphic monic polynomial families on $U$, of fixed degrees, and
are pairwise coprime for every $x\in U$. Let $P-P_0$ have degree less
than $n$ in $z$ and positive Hahn support with ordinary holomorphic
coefficients on $U$. Then \eqref{eq:henselfactor} has unique monic
Hahn-coherent factors reducing to the $p_i$. Their corrections have
support in the positive monoid generated by the input support, and
every coefficient is holomorphic on the original domain $U$.
\end{theorem}
\begin{proof}
The matrix $L(x)$ of \eqref{eq:factorlinear} is holomorphic on $U$.
It is invertible at every $x$ by pairwise coprimality, so
$L(x)^{-1}=\operatorname{adj}(L(x))/\det L(x)$ is holomorphic on all
of $U$. Apply \eqref{eq:henselrecursion} coefficientwise with this fixed
holomorphic matrix inverse. The same finite-word support argument
works uniformly, and no step requires a smaller domain. Uniqueness
again follows at the least differing exponent.
\end{proof}
This is a finite-dimensional parameter version of the matrix inversion
principle in \cite{SourceResearch}. It requires an ordinary coprime
factorization on the given base; it does not manufacture globally
labeled roots through ordinary monodromy or across a discriminant locus.

\subsection{An explicit two-factor recursion}
\begin{example}\label{ex:hensel}
For $P=z^2+tz-1+t^2$, the ordinary factors are $z-1$ and $z+1$.
The normalized roots have expansions
\[
 \alpha_+=1-\tfrac12t-\tfrac38t^2-\tfrac9{128}t^4+O(t^6),\qquad
 \alpha_-=-1-\tfrac12t+\tfrac38t^2+\tfrac9{128}t^4+O(t^6).
\]
These follow either from \eqref{eq:henselrecursion} or from
$\alpha_\pm=(-t\pm\sqrt{4-3t^2})/2$ and the strongly summable binomial
series. The two factors are $z-\alpha_+$ and $z-\alpha_-$. Their sum
and product have the exact prescribed coefficients; all displayed
support is in the positive integer monoid generated by the input.
\end{example}

\section{Optimal matching of perturbed root multisets}\label{sec:holder}

\subsection{Why a nearest-root estimate is not enough}
If $P,Q\in\OO_K[z]$ are monic of degree $n$ and their coefficients
agree to valuation at least $\varepsilon>0$, evaluating $Q-P$ at a
root of $P$ shows that some root of $Q$ is close. That alone does
not produce a bijection of root multisets: several roots could be
sent to the same nearby root. The next proof balances the multiplicities
in every relevant ultrametric cluster before matching.

\begin{theorem}[Optimal valuation-H\"older root matching]\label{thm:holder}
Let
\[
 P=z^n+\sum_{j<n}a_jz^j,\qquad
 Q=z^n+\sum_{j<n}b_jz^j,\qquad a_j,b_j\in\OO_K,
\]
where $n\ge1$. Suppose $v(a_j-b_j)\ge\varepsilon>0$ for every $j<n$.
There are listings of their roots, repeated with multiplicity, such that
\begin{equation}\label{eq:holder}
 \boxed{\displaystyle v(\alpha_i-\beta_i)\ge\frac{\varepsilon}{n}
 \qquad(1\le i\le n).}
\end{equation}
The exponent $\varepsilon/n$ is optimal uniformly over all such
polynomials.
\end{theorem}
\begin{proof}
All roots of both polynomials are finite. Put $\rho=\varepsilon/n$.
In the finite combined root multiset, inspect the valuations of the
nonzero pairwise differences that are strictly less than $\rho$.
Choose $\mu$ with $0<\mu<\rho$ and greater than all of these finitely
many valuations. Such a $\mu$ exists by divisibility: take the midpoint
between $\rho$ and the maximum of that finite set together with zero.

On the finite collection of root locations, the equivalence relations
$v(x-y)\ge\mu$ and $v(x-y)\ge\rho$ now coincide. Their classes are
intersections with closed $\mu$-balls. Take any class and choose a
root $a$ in it, from either polynomial. Translation by the finite
center $a$ preserves the lower valuation bound $\varepsilon$ for
coefficients of $E=Q-P$, since each translated coefficient is a finite
sum of an original coefficient times an integer and a nonnegative
power of $a$. Dilation by $t^\mu$, with $\mu>0$, gives
\[
 w_{a,\mu}(E)\ge\varepsilon.
\]
On the other hand, the monic root-product formula yields
\[
 w_{a,\mu}(P)=\sum_{i=1}^n\min\{\mu,v(\alpha_i-a)\}
 \le n\mu<\varepsilon.
\]
Thus \cref{cor:valrouche} says that this class contains exactly the
same number of $P$ roots and $Q$ roots, with multiplicity. This
argument also rules out a class containing roots of only one polynomial.
Match the two finite multisets arbitrarily within each class. Its
members have pairwise differences of valuation at least $\rho$, so
all these pairings satisfy \eqref{eq:holder}.

For optimality, take $P=z^n$ and $Q=z^n-t^\varepsilon$. Every root of
$Q$ has valuation exactly $\varepsilon/n$, whereas every root of $P$
is zero. No larger uniform bound is possible.
\end{proof}
The proof uses finitely many roots and their finite separation data.
There is no passage to a limit of roots as an auxiliary radius increases.
In particular it works when $\Gamma$ has arbitrary ordered-group rank.
The correspondence is not generally unique, as repeated or unresolved
clusters naturally allow multiple labelings.

\begin{corollary}[Simultaneous stability for derivatives]\label{cor:derivativeholder}
Under the hypotheses of \cref{thm:holder}, for every integer $0\le r<n$
the root multisets of $P^{(r)}$ and $Q^{(r)}$ admit a matching with
\[
 v(\alpha_i^{(r)}-\beta_i^{(r)})\ge\frac{\varepsilon}{n-r}.
\]
\end{corollary}
\begin{proof}
Divide each $r$th derivative by the nonzero integer
$n(n-1)\cdots(n-r+1)$ to make it monic of degree $n-r$.
Its coefficients remain finite, and the coefficient differences retain
valuation at least $\varepsilon$, because all nonzero rational factors
have valuation zero. Apply \cref{thm:holder}.
\end{proof}
For first derivatives the denominator is $n-1$, not $n$. It is sharp:
$P=z^n$ and $Q=z^n-t^\varepsilon z$ make the normalized derivative
change from $z^{n-1}$ to $z^{n-1}-t^\varepsilon/n$.

\subsection{What the integral normalization excludes}
The condition that the coefficients are finite is used twice: roots
are finite, and translating an error by a root does not reduce its
coefficient valuation bound. For arbitrary coefficients, first apply
\cref{cor:normalization} to a common root scale, then measure the
rescaled coefficient error. An unnormalized claim with the same
$\varepsilon/n$ at all infinite root locations would discard the
conditioning factors introduced by that change of variables.

\section{Separated roots, discriminants, and a Newton correction}
\label{sec:separated}

\subsection{A rootwise condition number}
Let $P\in\OO_K[z]$ be monic of degree $n\ge2$ with distinct roots
$\alpha_1,\ldots,\alpha_n$. Put
\begin{equation}\label{eq:di}
 d_i=v(P'(\alpha_i))=\sum_{j\ne i}v(\alpha_i-\alpha_j),\qquad
 \delta_i=\max_{j\ne i}v(\alpha_i-\alpha_j).
\end{equation}
Both numbers are nonnegative. The sum $d_i$ measures the size of the
inverse derivative at the root, while $\delta_i$ is the depth of its
nearest other root. Distinguishing them gives a sharper threshold than
a single coarse discriminant bound.

\begin{theorem}[Simple-root stability with a second-order error]
\label{thm:separated}
Let $Q=P+E$ be monic of degree $n$, where $\deg E<n$ and all
coefficients of $E$ have valuation at least $\varepsilon$. Suppose
\begin{equation}\label{eq:separationcondition}
 \varepsilon>\max_i(d_i+\delta_i).
\end{equation}
Then $Q$ has one simple root $\beta_i$ in each disjoint open ball
$\Bgt{\alpha_i}{\delta_i}$, and these are all its roots. For every $i$,
\begin{equation}\label{eq:rootwisebound}
 v(\beta_i-\alpha_i)\ge\varepsilon-d_i,
\end{equation}
and
\begin{equation}\label{eq:newtonerror}
 v\left(\beta_i-\alpha_i+
             \frac{E(\alpha_i)}{P'(\alpha_i)}\right)
 \ge2(\varepsilon-d_i)-\delta_i.
\end{equation}
If $E(\alpha_i)\ne0$, the exact displacement valuation is
\begin{equation}\label{eq:exactdisplacement}
 v(\beta_i-\alpha_i)=v(E(\alpha_i))-d_i.
\end{equation}
If $E(\alpha_i)=0$, then $\beta_i=\alpha_i$.
\end{theorem}
\begin{proof}
At center $\alpha_i$ and scale $\delta_i$, the product formula gives
\[
 w_{\alpha_i,\delta_i}(P)=\delta_i+d_i,
\]
since all the other root separations have valuation at most $\delta_i$.
The initial polynomial has a simple zero at the origin: precisely one
root lies in $\Bgt{\alpha_i}{\delta_i}$. Translation by $\alpha_i$
and nonnegative dilation preserve the error bound, so
$w_{\alpha_i,\delta_i}(E)\ge\varepsilon>d_i+\delta_i$.
Valuation Rouch\'e supplies exactly one $Q$ root in that open ball,
counted with multiplicity. It is therefore simple. Two of the indicated
balls cannot intersect, because their two centers have difference
valuation at most each of their defining depths. Their $n$ simple
roots exhaust the degree of $Q$.

Write $h=\beta_i-\alpha_i$. If $h=0$ the nontrivial estimates are
immediate. Otherwise $v(h)>\delta_i$, and the exact product expansion is
\begin{equation}\label{eq:productnearroot}
 P(\alpha_i+h)=hP'(\alpha_i)
       \prod_{j\ne i}\left(1+\frac{h}{\alpha_i-\alpha_j}\right).
\end{equation}
Every factor in parentheses is a unit with residue one. Hence
$v(P(\beta_i))=v(h)+d_i$. But
$P(\beta_i)=-E(\beta_i)$, and $\beta_i$ is finite, so
$v(E(\beta_i))\ge\varepsilon$. This proves \eqref{eq:rootwisebound}.

Let $R_P=P(\alpha_i+h)-P'(\alpha_i)h$ and
$R_E=E(\alpha_i+h)-E(\alpha_i)$. Expanding the finite product in
\eqref{eq:productnearroot} and the finite Taylor polynomial of $E$
gives
\begin{equation}\label{eq:remainderbounds}
 v(R_P)\ge d_i+2v(h)-\delta_i,\qquad
 v(R_E)\ge\varepsilon+v(h).
\end{equation}
The first inequality follows because every term of the product minus
one contains at least one $h/(\alpha_i-\alpha_j)$, of valuation at
least $v(h)-\delta_i$. For the second, every translated error
coefficient has valuation at least $\varepsilon$, and $h$ has positive
valuation.

The equation $Q(\alpha_i+h)=0$ becomes
\[
 P'(\alpha_i)h+E(\alpha_i)+R_P+R_E=0.
\]
Divide by $P'(\alpha_i)$, use \eqref{eq:remainderbounds}, and then
\eqref{eq:rootwisebound}. The resulting bound is
\[
 \min\{2v(h)-\delta_i,\ \varepsilon+v(h)-d_i\}
 \ge2(\varepsilon-d_i)-\delta_i,
\]
which proves \eqref{eq:newtonerror}.
Finally both remainders have valuation strictly greater than
$d_i+v(h)$, because $v(h)>\delta_i$ and $\varepsilon>d_i$.
The two terms of least valuation in the equation must therefore be
$P'(\alpha_i)h$ and $E(\alpha_i)$, proving
\eqref{eq:exactdisplacement}. If $E(\alpha_i)=0$, the original root
itself lies in the ball and uniqueness gives $\beta_i=\alpha_i$.
\end{proof}
The right side of \eqref{eq:newtonerror} is strictly greater than
$\varepsilon-d_i$ under \eqref{eq:separationcondition}. Thus
$-E(\alpha_i)/P'(\alpha_i)$ is a genuine first-order correction with
a strictly smaller error, not merely a formal suggestion to apply
Newton's method. No infinite Newton iteration is required for the proof.
For degree one the exact root displacement is simply the negative
constant coefficient error.

\subsection{Discriminant precision and a sharp failure at equality}
For a monic polynomial set
\[
 \Disc(P)=\prod_{i<j}(\alpha_i-\alpha_j)^2.
\]
This is a polynomial in its coefficients; the determinant interpretation
will be proved in \cref{sec:residue}. Its valuation for squarefree integral
$P$ is
\begin{equation}\label{eq:discvaluation}
 v(\Disc(P))=2\sum_{i<j}v(\alpha_i-\alpha_j)=\sum_i d_i.
\end{equation}

\begin{corollary}[A coefficient-level sufficient threshold]
\label{cor:discthreshold}
The stronger condition $\varepsilon>v(\Disc(P))$ implies every
conclusion of \cref{thm:separated}.
\end{corollary}
\begin{proof}
For each $i$, choose $j\ne i$ attaining $\delta_i$. Then
$d_j\ge\delta_i$, and all the other $d_\ell$ are nonnegative.
Consequently $d_i+\delta_i\le\sum_\ell d_\ell=v(\Disc(P))$.
\end{proof}
The discriminant gives a convenient global threshold, but the rootwise
maximum in \eqref{eq:separationcondition} can be substantially smaller.

\begin{example}[A collision exactly at the threshold]\label{ex:threshold}
For $a>0$ in the exponent group, let
\[
 P(z)=z(z-t^a),\qquad
 Q(z)=P(z)+\frac{t^{2a}}4=(z-t^a/2)^2.
\]
Both roots of $P$ have $d_i=\delta_i=a$. The coefficient error has
valuation $2a=d_i+\delta_i=v(\Disc(P))$, but $Q$ has a double root.
That root is not in either of the separated open balls
$\Bgt0a$ and $\Bgt{t^a}a$. Thus strictness in the general threshold
cannot be replaced by a weak inequality. This also realizes equality
in the global H\"older bound for $n=2$.
\end{example}

\subsection{Stability of the root tree and its critical directions}
\begin{corollary}[Preservation of pairwise scale data]\label{cor:treestability}
Under \cref{thm:separated}, the matching preserves all pairwise
separation valuations and their leading coefficients:
\[
 v(\beta_i-\beta_j)=v(\alpha_i-\alpha_j),\qquad
 \lc(\beta_i-\beta_j)=\lc(\alpha_i-\alpha_j).
\]
At every internal node of the original root-cluster tree, using its
original center and scale, $P$ and $Q$ have the same initial polynomial.
Their critical multiplicities in every residue direction of that node
therefore agree as well.
\end{corollary}
\begin{proof}
The two displacement errors have valuation strictly greater than
$\delta_i$ and $\delta_j$, hence greater than
$v(\alpha_i-\alpha_j)$. Adding them cannot change the least exponent
or leading coefficient of the difference.

A node may be centered at one of its roots $\alpha_i$, and its depth
$\rho$ is at most $\delta_i$. All roots are finite, so $\rho\ge0$.
Equation \eqref{eq:profileproduct} gives
\[
 w_{\alpha_i,\rho}(P)=\rho+\sum_{j\ne i}
                  \min\{\rho,v(\alpha_i-\alpha_j)\}
 \le\delta_i+d_i<\varepsilon.
\]
The weighted error is at least $\varepsilon$, so
\cref{cor:valrouche} gives equality of the initial polynomials.
At an internal node this polynomial is nonconstant, and
\cref{lem:initialderivative} gives equality for the derivatives too.
The direction counts now follow from \cref{thm:initialroots}.
\end{proof}
This preserves the root-scale hierarchy and the allocated critical
directions, not the individual multiplicities of all critical points.
For instance $z^3-1$ is squarefree but its derivative has a double
root at zero. Adding $t^Nz$, with $N>0$, leaves the ordinary root
clusters unchanged and splits that critical point into two at a
smaller scale. The direction count at the original scale is still two.

\begin{example}[Precision for the nested quartic]\label{ex:quarticperturb}
For the quartic of \cref{ex:nested}, in the root order $0,t^3,t,1$,
\[
 (d_i)=(4,4,2,0),\qquad (\delta_i)=(3,3,1,0).
\]
Thus the rootwise sufficient threshold is $\varepsilon>7$, whereas
the discriminant condition gives the more conservative $\varepsilon>10$.
For $Q=P+t^8(1+z)$, the four matched displacement valuations are exactly
$4,4,6,8$, by \eqref{eq:exactdisplacement}. Their first terms are
\[
 \beta_1=t^4+O(t^5),\quad
 \beta_2=t^3-t^4+O(t^5),\quad
 \beta_3=t+t^6+O(t^7),\quad
 \beta_4=1-2t^8+O(t^9).
\]
These follow from $E(\alpha_i)$ and the leading coefficients of
$P'(\alpha_i)$ at the four roots.
All three original cluster scales and the residual critical directions
$3/4,2/3,1/2$ survive this perturbation.
\end{example}

\section{Residues, resultants, and duality through collisions}\label{sec:residue}

\subsection{A functional requiring no contour}
Let $P\in K[z]$ be monic of degree $n\ge1$, and put $A_P=K[z]/(P)$.
For a polynomial $H$, let $R_H$ be its remainder on division by $P$.
Define
\begin{equation}\label{eq:lambda}
 \lambda_P(H)=[z^{n-1}]R_H.
\end{equation}
It is a well-defined $K$-linear functional on $A_P$. Since division by
a monic polynomial uses no inverse coefficients, if $P,H\in\OO_K[z]$
then $R_H$ and $\lambda_P(H)$ are integral as well.

\begin{theorem}[Collision-stable polynomial residue pairing]
\label{thm:residuepairing}
The bilinear pairing
\[
 A_P\times A_P\longrightarrow K,\qquad (G,H)\longmapsto\lambda_P(GH)
\]
is perfect, including when $P$ has repeated roots. In the basis
$1,z,\ldots,z^{n-1}$ its Gram determinant is
\begin{equation}\label{eq:resdet}
 \det G_P=(-1)^{n(n-1)/2}.
\end{equation}
The same pairing is perfect over $\OO_K$ on $\OO_K[z]/(P)$ when
$P$ is integral and monic.
\end{theorem}
\begin{proof}
If $i+j<n-1$, then $\lambda_P(z^{i+j})=0$, while for $i+j=n-1$ it
is one. Reverse the column order of the Gram matrix. Its entry in
row $i$ and column $j$ is then zero for $i<j$ and one for $i=j$.
Thus the reversed matrix is lower triangular with diagonal ones.
The column-reversal sign is $(-1)^{n(n-1)/2}$, proving
\eqref{eq:resdet}. This determinant is a unit over $K$ and over
$\OO_K$, so the pairing identifies the free rank-$n$ module with its
dual in either case.
\end{proof}
This finite algebraic result is the polynomial core of the
one-variable residue duality in \cite{SourceResearch}; it is not
claimed as a new general duality theorem. Its strength for the current
problem is that no discriminant denominator appears.

\subsection{Formal residues at distinct or colliding roots}
For a rational function, a residue at $\xi$ means the coefficient of
$(z-\xi)^{-1}$ in its formal Laurent expansion. The residue at infinity
is defined in the coordinate $u=1/z$ for the differential, so it is
the negative of the coefficient of $z^{-1}$ in the expansion of the
rational function at infinity.

\begin{proposition}[Finite residue formula]\label{prop:localresidue}
If $P=\prod_\xi(z-\xi)^{m_\xi}$, then for every polynomial $H$,
\begin{equation}\label{eq:residuesum}
 \lambda_P(H)=\sum_\xi\Res_{z=\xi}\frac{H(z)}{P(z)}\,dz.
\end{equation}
Writing $U_\xi(z)=P(z)/(z-\xi)^{m_\xi}$, its local summand is
\begin{equation}\label{eq:resjet}
 \frac1{(m_\xi-1)!}
 \left.\frac{d^{m_\xi-1}}{dz^{m_\xi-1}}
       \left(\frac{H(z)}{U_\xi(z)}\right)\right|_{z=\xi}.
\end{equation}
For squarefree $P$ this reduces to
\begin{equation}\label{eq:lagrangeresidue}
 \lambda_P(H)=\sum_{i=1}^n\frac{H(\alpha_i)}{P'(\alpha_i)}.
\end{equation}
\end{proposition}
\begin{proof}
Polynomial division gives $H/P=S+R_H/P$. The coefficient of $z^{-1}$
at infinity is the top coefficient of $R_H$, since $P$ is monic and
$\deg R_H<n$. A partial-fraction expansion of $R_H/P$ shows that
this coefficient is the sum of its simple-pole coefficients: terms
$(z-\xi)^{-j}$ with $j\ge2$ have no $z^{-1}$ coefficient at infinity.
This proves \eqref{eq:residuesum}. Locally $U_\xi(\xi)\ne0$, so its
inverse exists as a formal power series at $\xi$; extracting the
coefficient of degree $m_\xi-1$ in $H/U_\xi$ gives
\eqref{eq:resjet}. For $m_\xi=1$, $U_\xi(\xi)=P'(\xi)$.
\end{proof}
For example, when $P=z(z-t)$,
$\lambda_P(H)=(H(t)-H(0))/t$. The individual evaluation weights are
infinite, but for $H\in\OO_K[z]$ their sum is integral by
\eqref{eq:lambda}. At the repeated polynomial $P=(z-t)^2$, one has
$\lambda_P(H)=H'(t)$. This is an exact formula involving a jet, not
a limit taken as two roots approach each other.

The same partial-fraction argument gives the algebraic residue theorem
for every rational differential on the projective line: the sum of
its finitely many finite residues and its residue at infinity is zero.
Again this does not require a topological loop in $\SC$.

\subsection{Trace and norm}
For $H\in A_P$, let $m_H$ denote multiplication by $H$ on the
$n$-dimensional vector space $A_P$.

\begin{theorem}[Trace, norm, and the derivative trace element]
\label{thm:traceresidue}
For the distinct roots $\xi$ with multiplicities $m_\xi$,
\begin{align}
 \Tr(m_H)&=\sum_\xi m_\xi H(\xi)
          =\lambda_P(P'H),\label{eq:trace}\\
 \Nm(m_H):=\det(m_H)&=\prod_\xi H(\xi)^{m_\xi}.\label{eq:norm}
\end{align}
In particular, the element corresponding to the trace functional under
the perfect residue pairing is the class of $P'$.
\end{theorem}
\begin{proof}
In the local factor $K[u]/u^{m_\xi}$ from \cref{thm:crt}, multiplication
by $H$ is triangular in the powers of $u$, with every diagonal entry
$H(\xi)$. Its trace is $m_\xi H(\xi)$ and determinant
$H(\xi)^{m_\xi}$. Summing traces and multiplying determinants over
the finite product proves the first expression and \eqref{eq:norm}.

For the derivative expression, the coefficient of $z^{-1}$ at infinity
in $HP'/P$ equals $\lambda_P(P'H)$ by the same monic-division argument
as before. But \eqref{eq:logderivative} gives
$HP'/P=\sum_\xi m_\xi H(z)/(z-\xi)$, and the coefficient of $z^{-1}$
in $H(z)/(z-\xi)$ is $H(\xi)$, by finite polynomial division at
$\xi$. This proves \eqref{eq:trace}. Uniqueness of the representing
element follows from the perfect pairing.
\end{proof}
The trace cannot detect nilpotent directions in a local factor,
whereas the residue functional does detect the higher jets needed
for a perfect pairing. These are different bilinear forms and should
not be interchanged.

\subsection{Resultants and the discriminant determinant}
For monic $P$ of degree $n$ and arbitrary $Q$, define
\[
 \Res(P,Q)=\det(m_Q:A_P\to A_P).
\]
By \eqref{eq:norm},
\begin{equation}\label{eq:resultantproduct}
 \Res(P,Q)=\prod_{i=1}^nQ(\alpha_i).
\end{equation}
For $Q\ne0$ of degree $m$ and nonmonic $P$ with leading coefficient
$a_n$, the usual normalization is
$\Res(P,Q)=a_n^m\prod_iQ(\alpha_i)$. This agrees with the Sylvester
determinant convention in which $\Res(z-\alpha,Q)=Q(\alpha)$.
One way to check the determinant identity is to use division by the
monic $P$ in the map $(U,V)\mapsto UP+VQ$ with
$\deg U<m$, $\deg V<n$: its quotient block is an identity and its
remainder block is multiplication by $Q$ modulo $P$. The coefficient
bases give the conventional Sylvester sign and the factor $a_n^m$
after normalization.

Thus the resultant vanishes exactly when $P,Q$ have a common root.
Equivalently, it vanishes when the B\'ezout map is singular. This
provides a coefficient-level test, but does not promise an effective
algorithm for equality of arbitrary infinitely supported surreal
coefficients.

\begin{theorem}[Discriminant and trace-pairing degeneration]
\label{thm:discriminant}
For monic $P$ of degree $n$, let
$T_{ij}=\Tr(m_{z^{i+j}})$, $0\le i,j<n$. Then
\begin{equation}\label{eq:tracedisc}
 \det T=\Disc(P)
 =(-1)^{n(n-1)/2}\Res(P,P').
\end{equation}
The trace pairing is nondegenerate exactly when $P$ is squarefree,
whereas the residue pairing in \cref{thm:residuepairing} is always
nondegenerate.
\end{theorem}
\begin{proof}
List the $n$ roots with repetitions, and let $V_{ki}=\alpha_k^i$.
The trace formula says $T=V^TV$, where transpose, not conjugate
transpose, is used. Hence
\[
 \det T=(\det V)^2=\prod_{i<j}(\alpha_i-\alpha_j)^2.
\]
This identity remains valid with repeated rows, so it handles repeated
roots directly. Also
$\prod_iP'(\alpha_i)=(-1)^{n(n-1)/2}
\prod_{i<j}(\alpha_i-\alpha_j)^2$: in the squarefree case multiply
$P'(\alpha_i)=\prod_{j\ne i}(\alpha_i-\alpha_j)$, and with repeated
roots both sides vanish. Apply \eqref{eq:resultantproduct}.
The product of squared differences is nonzero precisely for distinct
roots. The residue determinant is the fixed unit
\eqref{eq:resdet}, proving the contrast.
\end{proof}
This determinant formula proves that the earlier product definition of
the discriminant is a coefficient polynomial, and links its valuation
to the root-separation precision in \cref{sec:separated}.

\subsection{An explicit cubic duality calculation}
The following calculation also appears in the supplied several-variable
manuscript \cite{SourceResearch}; it is included here as a normalization
check and a bridge to that work, not as a new example.
Let $P=z^3-az-b$. In the basis $(1,z,z^2)$,
\[
 G_P=\begin{pmatrix}0&0&1\\0&1&0\\1&0&a\end{pmatrix},\qquad
 T_P=\begin{pmatrix}3&0&2a\\0&2a&3b\\2a&3b&2a^2\end{pmatrix}.
\]
Their determinants are $-1$ and $4a^3-27b^2$, respectively. The
residue-dual basis is $(z^2-a,z,1)$, so its trace element is
$(z^2-a)+z^2+z^2=3z^2-a=P'$.

For $a=3t^2$, $b=2t^3$ the polynomial is
$(z-2t)(z+t)^2$. The residue and trace are
\[
 \lambda_P(H)=\frac{H(2t)-H(-t)}{9t^2}-\frac{H'(-t)}{3t},\qquad
 \Tr(m_H)=H(2t)+2H(-t).
\]
The first identity follows by partial fractions or \eqref{eq:resjet}.
The derivative term retains the nilpotent information at the double
point that the trace omits.

\section{Several polynomial variables in a set-sized workspace}\label{sec:multivariate}
The one-variable results fit into ordinary algebraic geometry over $K$.
It is useful to make the legitimate transfer explicit, because the
proper-class notation $\SC$ by itself is not a substitute for a
set-sized ring in a scheme construction.

\subsection{Polynomial equations and the Nullstellensatz}
\begin{lemma}[The finite-generation field lemma]\label{lem:zariski}
If a field $L$ is finitely generated as a $K$-algebra, with $K$
algebraically closed, then $L=K$.
\end{lemma}
\begin{proof}
Among a finite algebra generating list choose a maximal algebraically
independent subset $t_1,\ldots,t_r$. The remaining generators are
algebraic over $K(t_1,\ldots,t_r)$. Clearing the finitely many
coefficient denominators of their monic algebraic equations supplies
a nonzero $f\in K[t_1,\ldots,t_r]$ such that all the remaining
generators are integral over
$B=K[t_1,\ldots,t_r,f^{-1}]$. Because $L$ is a field containing the
original generators, $L$ equals the $B$-algebra generated by those
remaining elements and is integral over $B$.

An integral subring of a field over which the whole field is integral
must itself be a field: for $0\ne b\in B$, a monic equation for
$b^{-1}$, multiplied by a suitable power of $b$, expresses $b^{-1}$
as an element of $B$. If $r\ge1$, however, choose $c\in K$ such that
$t_1-c$ does not divide $f$; only finitely many $c$ are excluded.
Then $t_1-c$ is not a unit in $B$, since its quotient ring after
localization is a nonzero localization of $K[t_2,\ldots,t_r]$.
This contradicts $B$ being a field. Thus $r=0$, so $L$ is a finite
algebraic extension of the algebraically closed field $K$, and $L=K$.
\end{proof}

\begin{theorem}[Surcomplex finite-system Nullstellensatz]
\label{thm:nullstellensatz}
Let $I\subset K[x_1,\ldots,x_d]$ be an ideal in a set-sized
algebraically closed Hahn workspace. It is proper if and only if
its polynomials have a common zero in $K^d$. Moreover a polynomial
$H$ vanishes at all common zeros of $I$ if and only if some positive
power of $H$ belongs to $I$.

Consequently a finite polynomial system with surcomplex coefficients
has a surcomplex solution exactly when its generated ideal is proper
in any algebraically closed workspace containing those coefficients;
when solvable, it already has a solution in that workspace.
\end{theorem}
\begin{proof}
A proper ideal lies in a maximal ideal $\mathfrak n$. Its quotient
is a field finitely generated as a $K$-algebra, hence equals $K$ by
\cref{lem:zariski}. The images of the coordinates give the common
zero. Conversely evaluation at a common zero is a homomorphism
annihilating $I$, so $1\notin I$.

If $H$ vanishes at all common zeros, consider the ideal
$(I,1-yH)$ in one additional variable. A common zero would have
$H=0$ and $yH=1$, impossible. By the first assertion this ideal is
the whole ring. Express $1$ as a finite combination of its generators,
substitute $y=H^{-1}$ in the localization, and multiply by a high
enough power $H^N$ to clear denominators. This gives $H^N\in I$.
The converse is immediate by evaluation. The zero polynomial case is
trivial. This is the usual Rabinowitsch argument; compare the standard
Nullstellensatz formulation \cite{Stacks}.

For the final assertion, a proper workspace ideal has a solution in
that workspace by what was just proved. If it is the whole ring,
a finite identity $1=\sum A_if_i$ persists in every extension, so
there can be no surcomplex solution.
\end{proof}
No priority claim attaches to this classical theorem. The point is
that the field and ideal to which it is applied are specified and
set-sized. A positive-dimensional solution set can acquire new points
in a larger workspace, even though solvability does not change.

\subsection{Finite quotient algebras retain actual multiplicities}
\begin{proposition}[Finite fibers and local lengths]\label{prop:finitealgebra}
If $A=K[x_1,\ldots,x_d]/I$ is finite-dimensional of dimension $N>0$,
then its common-zero set $Z\subset K^d$ is finite and
\[
 A\cong\prod_{\xi\in Z}A_\xi,
 \qquad
 \sum_{\xi\in Z}\dim_K A_\xi=N.
\]
Each $A_\xi$ is local with residue field $K$ and nilpotent maximal
ideal. Thus $|Z|\le N$, with equality precisely when all local factors
have dimension one. After extension to a larger algebraically closed
workspace there are no new point locations, and the local dimensions
remain unchanged.
\end{proposition}
\begin{proof}
Distinct maximal ideals of $A$ are comaximal. By the finite Chinese
remainder theorem, any $s$ distinct ones give a surjection $A\to K^s$,
using \cref{lem:zariski} for their residue fields. Hence there are at
most $N$ maximal ideals. They correspond to the points $Z$ by
\cref{thm:nullstellensatz}. Let $J$ be their intersection, the
Jacobson radical. The descending chain $J^m$ stabilizes because $A$
is finite-dimensional. For a stable term $M=J^m=JM$, the determinant
argument proves $M=0$: choose finitely many generators of $M$, write
each as a linear combination of them with coefficients in $J$, and
multiply the resulting matrix equation by its adjugate. Its determinant
is of the form $1+j$ with $j\in J$, hence is a unit, and annihilates
all generators. Therefore $J^m=0$ for some $m$.

The powers $\mathfrak n_\xi^m$ of the finitely many maximal ideals
are pairwise comaximal and have intersection
$\prod_\xi\mathfrak n_\xi^m=J^m=0$. The Chinese remainder theorem
now gives $A\cong\prod_\xi A/\mathfrak n_\xi^m$. Each factor has
one maximal ideal, which is nilpotent, and residue field $K$.
Dimensions add. After scalar extension, each such factor still has
a nilpotent ideal and quotient the extension field; it is still local
and of the same vector-space dimension. The coordinates modulo that
ideal are the original tuple $\xi$, so no new point is introduced.
\end{proof}
The integer $\dim_K A_\xi$ is the algebraic local multiplicity.
This is the finite polynomial counterpart of the actual-fiber
identification in \cite{SourceResearch}. For an arbitrary analytic
system, constructing such a finite algebra is additional work; it is
not supplied by this proposition alone.

\subsection{A coupled infinitesimal four-point fiber}
\begin{example}\label{ex:coupled}
Consider
\[
 f_1=x^2-ty-t^2,\qquad f_2=y^2-tx,
\]
with $t=\omega^{-1}$. Since $t$ is a nonzero element of the field,
set $Y=y/t$. The second equation gives $x=tY^2$, and the first becomes
$t^2(Y^4-Y-1)=0$. Hence
\begin{equation}\label{eq:couplediso}
 K[x,y]/(f_1,f_2)\cong K[Y]/(Y^4-Y-1).
\end{equation}
The quotient has dimension four. Its basis $(1,x,y,xy)$ maps to
the basis $(1,tY^2,tY,t^2Y^3)$ on the right.

The ordinary polynomial $Y^4-Y-1$ has four distinct roots. Indeed,
a common root with $4Y^3-1$ would satisfy $Y^3=1/4$ and
$Y/4-Y-1=0$, hence $Y=-4/3$, which does not have cube $1/4$.
Its roots $c_1,\ldots,c_4\in\C$ are all nonzero. Thus the actual
solutions are exactly
\[
 (x,y)=(tc_j^2,tc_j),\qquad 1\le j\le4.
\]
They are infinitesimal and simple: the Jacobian determinant
$4xy-t^2$ equals $t^2(4c_j^3-1)\ne0$ at each point.

In the stated basis the multiplication matrices are
\[
 W_x=\begin{pmatrix}
 0&t^2&0&0\\1&0&0&t^2\\0&t&0&t^2\\0&0&1&0
 \end{pmatrix},\qquad
 W_y=\begin{pmatrix}
 0&0&0&t^3\\0&0&t&0\\1&0&0&t^2\\0&1&0&0
 \end{pmatrix}.
\]
They commute and satisfy
$W_x^2=tW_y+t^2I$, $W_y^2=tW_x$.
Eliminating $x$ gives the resultant
$y^4-t^3y-t^4$. This example shows explicitly how a coupled system
whose ordinary reduction is $x^2=y^2=0$ can split into four actual
surcomplex points while its finite-algebra length stays four.
\end{example}

\section{Polynomials versus Hahn-coherent entire data}\label{sec:comparison}

\subsection{Uniform degree is the dividing line}
There is a simple exact description of the intersection of the two
languages. For fixed $\Gamma$,
\begin{equation}\label{eq:uniformdegree}
 K_\Gamma[z]
 =\bigcup_{d\ge0}\bigl(\C[z]_{\le d}((t^\Gamma))\bigr).
\end{equation}
The equality means equality of polynomial coefficient data, or of their
Hahn-coherent evaluations on the finite thickened plane.

To prove it, a finite-degree polynomial has finitely many Hahn
coefficients; their finite union is well ordered, and grouping by
exponent gives ordinary polynomial coefficient functions with the same
uniform degree bound. Conversely, a Hahn family of ordinary polynomials
of degree at most $d$ can be regrouped into $d+1$ Hahn coefficients.
All their supports are subsets of the given common well-ordered support,
so the regrouping produces an element of $K_\Gamma[z]$.

The bound must be uniform in the exponent. The coherent datum
\begin{equation}\label{eq:geometricentire}
 F(z)=\sum_{k\ge0}t^kz^k\in\OO(\C)((t^\Z))
\end{equation}
has ordinary entire coefficient functions and a well-ordered support,
but no uniform degree bound. On finite surcomplex points it equals
$1/(1-tz)$ by strong geometric summation. It is not a polynomial:
its ordinary coefficient functions have unbounded degrees, and
coefficient uniqueness excludes a bounded-degree representation.
Its rational continuation has a pole at the infinite point $1/t$.
Thus ``entire on the thickened ordinary plane'' is not ``entire at
all surreal scales.'' This elementary example is consistent with the
all-scale polynomial rigidity theorem stated in \cite{SourceAnalysis}.

\subsection{Finite support unions and the infinite global obstruction}
For a finite root multiset, any supports of its finitely many symmetric
coefficients have well-ordered union. Consequently the finite CRT,
cluster-factor, and residue constructions above never encounter a
failure caused merely by an infinite number of spatial centers.
They can still produce negative exponents through small denominators,
and root extraction can enlarge the exponent group.

The supplied global manuscript \cite{SourceGlobal} proves a different
phenomenon for infinitely many ordinary centers. Its cluster family
$P_n(u)=u-t^{1/n}$ cannot be the exact zero divisor of one Hahn-coherent
entire datum, because the union of the symmetric coefficient supports
is decreasing. In contrast its family $P_n(u)=u^n-t$ is realizable:
all symmetric coefficients use a single exponent, even though the
individual root supports contain the descending sequence $1/n$.
These are results of that supplied manuscript, not newly established
global theorems in the present article.

The polynomial lessons are compatible with this distinction.
Resultants, characteristic polynomials, remainders, and residue
pairings retain finite symmetric data even when root labels introduce
fractional scales or inverse separations. A claim that every separately
specified infinite family can be assembled from the finite results
would be unjustified.

\subsection{What has and has not been transferred}
Our polynomial Rouch\'e theorem quantifies over an order circle and
uses a fixed-degree first-order statement. Our valuation Rouch\'e
theorem instead compares finite coefficient valuations and preserves
an initial polynomial. The supplied contour theory integrates ordinary
holomorphic coefficients and then assembles a strongly summable Hahn
series. These are three different mechanisms. Agreement of some
consequences does not make the mechanisms interchangeable.

Similarly, the matrix proof of \cref{thm:schoenberg} uses the positive
definite form on the finite-dimensional space $K^n$ and algebraic
square roots. It does not assert completeness of an infinite-dimensional
surcomplex Hilbert space. The factor-lifting theorem uses a common
positive support, not a norm-convergent sequence of approximate
factorizations. The finite-system Nullstellensatz uses set-sized
polynomial rings, not a scheme whose underlying coordinate ring is
an unspecified proper class.

\section{A consolidated theorem map and mathematical boundaries}\label{sec:map}

\subsection{The classical analogues and the extra information proved}
\begin{longtable}{@{}P{0.26\textwidth}P{0.30\textwidth}P{0.35\textwidth}@{}}
\toprule
Classical principle & Surcomplex conclusion & Proof mechanism and status\\
\midrule
\endhead
Fundamental theorem of algebra & Finite factorization in a set-sized
Hahn workspace & Established Hahn algebraic closedness, then finite
induction; \cref{thm:fta}.\\[0.5em]
Gauss--Lucas and Jensen & Convex-hull and conjugate-disk bounds with
surreal-valued geometry & Direct finite reciprocal identities;
\cref{thm:gausslucas,thm:jensen}.\\[0.5em]
Schoenberg inequality & Sharp second-moment bound and collinearity
equality case & Finite matrix compression over a real-closed field;
\cref{thm:schoenberg}. Classical inequality, not a new conjecture
resolution.\\[0.5em]
Polynomial Rouch\'e & Order-circle counts at arbitrary surreal radii &
Precisely encoded real-closed-field transfer;
\cref{thm:ordrouche}.\\[0.5em]
Newton polygon & Counts in every ball, shell, and residue direction &
Finite initial-polynomial product; arbitrary ordered-group rank;
\cref{thm:initialroots,thm:newton}.\\[0.5em]
Non-Archimedean finite maps & Exact image balls and constant fiber
multiplicity & Apply the same product theorem to $P-y$;
\cref{thm:imageballs}.\\[0.5em]
Rolle/Gauss--Lucas at a scale & Exactly $k-1$ critical multiplicities
in an occupied $k$-root ball; residual branching directions &
Characteristic-zero differentiation of initial polynomials;
\cref{thm:criticalballs,thm:directions}.\\[0.5em]
Hensel factorization & Unique coprime factors on a specified common
support monoid & Finite CRT linearization and well-founded recursion;
\cref{thm:hensel}.\\[0.5em]
Continuity of root multisets & Optimal $\varepsilon/n$ valuation
matching, with multiplicity & Finite balanced valuation-ball partition;
\cref{thm:holder}.\\[0.5em]
Simple-root perturbation & Rootwise threshold, exact displacement
valuation, second-order Newton error & Product expansion plus
valuation Rouch\'e; \cref{thm:separated}.\\[0.5em]
Residue and trace duality & Perfect residue pairing through collisions;
trace determinant is the discriminant & Monic division, finite CRT,
formal Laurent coefficients; \cref{thm:residuepairing,thm:discriminant}.
Overlap with supplied results is credited.\\[0.5em]
Nullstellensatz and local multiplicity & Finite-system solvability and
finite local-algebra lengths & Classical algebra in each workspace;
\cref{thm:nullstellensatz,prop:finitealgebra}.\\
\bottomrule
\end{longtable}

\subsection{Hypotheses that cannot be discarded}
The degree is always finite. Without that assumption, neither the
finite factorization nor the finite root partitions and determinant
proofs are statements of the same kind. Divisibility of the exponent
workspace is available when roots or scale midpoints are needed;
no claim is made that $K_\Z$ is algebraically closed. Positive support
in \cref{thm:hensel} is a common well-ordered support condition, not
merely a pointwise assertion about an unrelated family of errors.

The exact critical count requires a ball occupied by at least one
root of the polynomial being counted; the more general mapping-ball
form uses $P-P(a)$ to ensure that occupation. All derivative arguments
use residue characteristic zero. Integral coefficients in the
perturbation estimates ensure that translating by a finite root
preserves the coefficient error bound. The sharp collision example
shows why the strict separation threshold cannot be suppressed.

The original root tree specifies an allocation of critical points and
certain residue directions, not their full Hahn expansions. When a
residual critical polynomial has a multiple root, that critical
subcluster may require its own rescaling analysis. No general
closed-form solution by radicals is asserted, and the presence of
surreal coefficients does not make ordinary higher-degree Galois
obstructions disappear in a prescribed smaller coefficient field.

\subsection{Literature assessment and further directions}
The foundational surreal and Hahn facts, the real-closed-field transfer
principle, the classical root geometry, and the general algebraic
geometry used here are established results. The source comparison
includes the primary work of Poonen, Pereira, Kushel--Tyaglov,
Eisermann, van den Dries, and the school-hosted non-Archimedean lecture
notes of Benedetto. The latter already provide disk mapping and
valuation-polygon methods in a complete rank-one analytic setting.
These precedents rule out presenting ``a non-Archimedean polynomial
theory'' or the Schoenberg inequality as a newly discovered general
principle.

The present contribution is a self-contained surcomplex development
that carries the stated hypotheses, arbitrary-rank scale bookkeeping,
root-multiset matching, residual critical directions, and support bounds
through explicit proofs. We do not assert that each precise estimate
or formulation is absent from all earlier valuation-theoretic
literature. The attached source manuscripts are treated as supplied
research documents, not as independently verified publications.

A natural continuation would develop an effective representation of
root-cluster refinements for restricted computable classes of surreal
coefficients, together with certified precision propagation for
resultants and multiple critical clusters. Another would compare
support-controlled parameter factors with ramification and monodromy
across the ordinary discriminant. These are directions for additional
work, not conclusions silently assumed in the present proofs.

\appendix
\section{A finite-dependency and support audit}\label{app:support}
The infinite operation in the main algebra is coprime factor lifting.
All root-product identities, derivatives, matrices, residue calculations,
and CRT decompositions involve only finitely many factors or coordinates.
For lifting, an output exponent $\nu$ receives finitely many support-word
contributions by \cref{lem:support}; this remains true for a group of
arbitrary rank and for well-ordered supports with infinitely many terms
below some upper bound.

\begin{center}
\begin{tabularx}{\textwidth}{@{}LL@{}}
\toprule
Operation & Size or support control\\
\midrule
Finite polynomial addition and multiplication & Finite unions and
Minkowski sums of coefficient supports.\\[0.4em]
Monic division of a finite polynomial & Finitely many field operations;
integral coefficients remain integral.\\[0.4em]
Initial polynomial at $(a,\rho)$ & One finite translation and dilation,
one minimum, and finitely many standard parts.\\[0.4em]
Newton profile and root-cluster tree & Finitely many active indices or
root pair separations; no cofinal sequence of scales.\\[0.4em]
Coprime factor correction & All outputs lie in $S^*$; at each exponent
only earlier correction coefficients occur.\\[0.4em]
Parameter factor correction & The same $S^*$ and a holomorphic finite
matrix inverse on the original ordinary domain.\\[0.4em]
H\"older matching & A finite partition of the combined root multiset;
no choice of an infinite limiting matching.\\[0.4em]
Separated-root Newton estimate & Exact finite product expansion with
explicit valuation remainders.\\[0.4em]
Residue and trace & Finite quotient matrices or finitely many formal
local Laurent coefficients; no contour limit.\\
\bottomrule
\end{tabularx}
\end{center}
A finite algebraic operation on arbitrary Hahn inputs is mathematically
well defined but not automatically computationally effective. A support
may have no supplied computable enumeration, and equality of arbitrary
surreal coefficients may have no finite decision procedure in the
chosen representation. The symbolic tests in the archive deliberately
use exact rational, polynomial, and finitely truncated formal-series
inputs, where the operations being tested are effective.

\section{Verification artifacts and reproducibility}\label{app:verification}
The archive includes \texttt{verify\_examples.py} and its output
\texttt{verification\_report.txt}. The script uses exact symbolic
arithmetic to check the following finite formulas: the cubic Newton
profiles and displayed approximate roots; the nested quartic initial
polynomials and critical directions; the coprime factor example; the
sharp quadratic collision; the nested-quartic separation data and
leading perturbations; the cubic residue and trace matrices; finite
instances of the compression identity and its squared-norm formula;
and the coupled two-variable multiplication matrices and resultant.
It also tests root-matching cluster counts for finitely many generated
rational-power polynomial examples.

These calculations are consistency checks, not substitutes for the
theorems. In particular, they do not certify arbitrary Hahn supports,
real-closed-field quantifier elimination, the universal factor-lifting
recursion, or arbitrary-degree root matching. The general claims rest
on the proofs in this article. No proof-assistant verification or
independent refereeing is asserted.

The LaTeX source has an internal bibliography and no external image or
bibliography dependencies. Compile with
\begin{verbatim}
latexmk -pdf surcomplex_polynomial_algebra.tex
\end{verbatim}
or run \texttt{pdflatex} repeatedly until references stabilize.
The Python checks require Python 3 and SymPy and are run by
\begin{verbatim}
python verify_examples.py
\end{verbatim}
The included report records the exact cases and truncation orders
actually checked. In the supplied run, all 189 exact checks passed,
including 20 generated matching cases of degrees two through six.
This finite test count is not a formal-verification claim.

\section{Notation index}\label{app:notation}
\begin{longtable}{@{}P{0.29\textwidth}P{0.63\textwidth}@{}}
\toprule
Notation & Meaning\\
\midrule
\endhead
$\SC=\No[i]$ & Surcomplex class field.\\[0.3em]
$t^\gamma=\omega^{-\gamma}$ & Conway normal-form monomial.\\[0.3em]
$F_\Gamma,K_\Gamma$ & Real and complex Hahn workspaces, with divisible
set-sized ordered exponent group $\Gamma$.\\[0.3em]
$|z|,v(z),\lc(z)$ & Surreal-valued modulus, natural valuation, and
leading complex coefficient.\\[0.3em]
$\OO_K,\mm_K,\st$ & Finite elements, infinitesimals, and standard part.\\[0.3em]
$D(a,r),\overline D(a,r)$ & Open and closed order disks defined by the
modulus.\\[0.3em]
$\Bge a\rho,\Bgt a\rho$ & Valuation balls defined by $v(z-a)\ge\rho$
or $v(z-a)>\rho$.\\[0.3em]
$w_{a,\rho}(P),I_{a,\rho}(P)$ & Weighted coefficient valuation and
normalized ordinary initial polynomial.\\[0.3em]
$j_-,j_+$ & Smallest and largest indices attaining the weighted
coefficient minimum.\\[0.3em]
$\phi_P$ & Newton valuation profile at center zero.\\[0.3em]
$E^*$ & Additive monoid generated by a positive well-ordered support,
including zero.\\[0.3em]
$d_i,\delta_i$ & Valuation of $P'(\alpha_i)$ and maximal separation
valuation from $\alpha_i$ to another root.\\[0.3em]
$A_P,\lambda_P$ & Finite quotient $K[z]/(P)$ and its top-remainder
coefficient functional.\\[0.3em]
$m_H$ & Multiplication by the class of $H$ on a finite algebra.\\[0.3em]
$\Res(P,Q),\Disc(P)$ & Polynomial resultant and discriminant.\\[0.3em]
$\Res_{z=\xi}R(z)\,dz$ & Formal local residue, a Laurent coefficient.\\[0.3em]
$\HH(U)$ & Hahn families of ordinary holomorphic functions on one
common ordinary domain $U$.\\
\bottomrule
\end{longtable}

\begin{thebibliography}{99}
\addcontentsline{toc}{section}{References}
\small
\bibitem{SourceAnalysis}
\emph{Surcomplex Analysis: Infinitesimal Calculus, Hahn-Coherent
Holomorphy, and Contour Theory}. Supplied manuscript dated September 21,
2026; file \texttt{surcomplex\_analysis(3).tex}. In particular, the
sections on workspaces, preparation, Rouch\'e, and all-scale rigidity.
No publication or authorship attribution is inferred from the filename.

\bibitem{SourceResearch}
\emph{Hartogs Coherence, Finite Maps, and Residue Duality over the
Surcomplex Numbers}. Supplied manuscript dated September 21, 2026;
file \texttt{surcomplex\_research.tex}. In particular, the sections on
positive operators, parameter division, finite fibers, and one-variable
residue and trace formulas.

\bibitem{SourceGlobal}
\emph{Global Divisors and Cohomological Obstructions in Hahn-Coherent
Surcomplex Analysis}. Supplied manuscript dated September 21, 2026;
file \texttt{surcomplex\_global\_theorems.tex}. In particular, the
symmetric-support divisor criterion and the support-restricted Chinese
remainder theorem.

\bibitem{Conway}
John H. Conway, \emph{On Numbers and Games}, second edition,
A K Peters, 2001; first edition, Academic Press, 1976.

\bibitem{Gonshor}
Harry Gonshor, \emph{An Introduction to the Theory of Surreal Numbers},
London Mathematical Society Lecture Note Series \textbf{110},
Cambridge University Press, 1986.

\bibitem{Neumann}
B. H. Neumann, ``On ordered division rings,''
\emph{Transactions of the American Mathematical Society}
\textbf{66} (1949), 202--252.

\bibitem{Poonen}
Bjorn Poonen, ``Maximally complete fields,''
\emph{L'Enseignement Math\'ematique} (2) \textbf{39} (1993), 87--106.
\href{https://math.mit.edu/~poonen/papers/amsval.pdf}{Author-hosted text}.
Corollary 4 supplies the divisible-group algebraic-closure input.

\bibitem{TarskiSurvey}
Lou van den Dries, ``Alfred Tarski's elimination theory for real closed
fields,'' \emph{The Journal of Symbolic Logic} \textbf{53} (1988),
no.~1, 7--19.
\href{https://doi.org/10.2307/2274424}{doi:10.2307/2274424}.

\bibitem{Eisermann}
Michael Eisermann, ``The fundamental theorem of algebra made effective:
an elementary real-algebraic proof via Sturm chains,''
\emph{The American Mathematical Monthly} \textbf{119} (2012), no.~9,
715--752.
\href{https://doi.org/10.4169/amer.math.monthly.119.09.715}{doi:10.4169/amer.math.monthly.119.09.715};
\href{https://arxiv.org/abs/0808.0097}{arXiv:0808.0097}.

\bibitem{Eberl}
Manuel Eberl, ``Two theorems about the geometry of the critical points of a
complex polynomial,'' \emph{Archive of Formal Proofs}, November 14, 2023.
\href{https://isa-afp.org/entries/Polynomial_Crit_Geometry.html}{Entry and formal development}.
Used as a reference for ordinary Gauss--Lucas and Jensen geometry,
not as a formal verification of this article.

\bibitem{Pereira}
Rajesh Pereira, ``Differentiators and the geometry of polynomials,''
\emph{Journal of Mathematical Analysis and Applications}
\textbf{285} (2003), no.~1, 336--348.
\href{https://doi.org/10.1016/S0022-247X(03)00465-7}{doi:10.1016/S0022-247X(03)00465-7}.

\bibitem{KushelTyaglov}
Olga Kushel and Mikhail Tyaglov, ``Circulants and critical points of
polynomials,'' \emph{Journal of Mathematical Analysis and Applications}
\textbf{439} (2016), no.~2, 634--650.
\href{https://doi.org/10.1016/j.jmaa.2016.03.005}{doi:10.1016/j.jmaa.2016.03.005};
\href{https://arxiv.org/abs/1512.07983}{arXiv:1512.07983}.
The Schoenberg inequality and the matrix-differentiator precedents
are classical inputs to the literature comparison.

\bibitem{Benedetto}
Robert L. Benedetto, \emph{Non-Archimedean Dynamics in Dimension One:
Lecture Notes}, Arizona Winter School, 2010.
\href{https://swc-math.github.io/aws/2010/2010BenedettoNotes-09Mar.pdf}{School-hosted notes}.
See \S3, especially Corollary 3.11 and the discussion of valuation
polygons, for the complete rank-one analytic setting.

\bibitem{Stacks}
The Stacks Project Authors, \emph{The Stacks Project},
\href{https://stacks.math.columbia.edu/tag/00FV}{Tag 00FV:
Hilbert Nullstellensatz}. Consulted September 21, 2026.

\bibitem{UserReference}
\emph{Surreal number}, subsection ``Surcomplex numbers,''
Wikipedia, the orientation reference specified in the request.
\href{https://en.wikipedia.org/wiki/Surreal_number\#Surcomplex_numbers}{Referenced subsection}.
The proofs use the primary and supplied sources identified above,
not this overview as an independent algebraic-closure proof.
\end{thebibliography}
\end{document}
